
\section{Anomaly-completed topological holography from first principles}
% label removed: app:anomaly-completed-topological-holography

\index{topological holography|textbf}
\index{Koszul duality}
\index{chiral algebra}
\index{factorization algebra}
\index{anomaly completion}
\index{transgression algebra}
\index{genus-Clifford completion}
\index{Yangian}
\index{affine Kac--Moody algebra}
\index{Fock module}

The surrounding chapters share a common algebraic structure:
bulk algebras are classical or shifted-Poisson/chiral, boundary algebras are their quantum chiral shadows, line operators are controlled by Koszul-dual algebras, and higher operations encode the deformation and anomaly data obstructing a naive strictification of the bulk--boundary dictionary.
The precise mathematical form is not an equality of algebras but a \emph{duality plus anomaly completion}.
This appendix isolates that algebraic core, corrects it at first principles, and pushes it one step farther.

The literature inputs are:
chiral algebras as quantum avatars of coisson structures \cite{BD04};
higher-dimensional chiral/factorization Koszul duality \cite{FG12};
derived centers and derived tensor products in topological field theory \cite{costello-gaiotto,CP2020,DNP25};
higher operations and Wess--Zumino consistency in holomorphic-topological perturbation theory \cite{GKW24};
boundary chiral algebras in holomorphic twists \cite{CDG20};
twisted holography via Koszul duality \cite{Zeng23};
and line operators controlled by Yangian-type structures \cite{DNP25}.
The theorem package below turns this constellation into a strict dg-algebraic machine.

\medskip

\noindent\textbf{What is proved here.}
The proved core of this appendix is the construction of:
\begin{enumerate}
\item a universal \emph{transgression algebra} \(B_{\Theta}\) attached to a dg algebra \(B\) and a closed central degree-\(2\) class \(\Theta\);
\item a corrected one-sided two-term resolution computing both \(R\!\operatorname{Hom}\) and derived relative tensor products in the anomaly-completed theory;
\item a \emph{secondary anomaly} \(u=\eta^{2}\), together with a quadratic refinement on the space of neutralizations;
\item a \emph{genus-Clifford completion} \(\mathfrak{G}_{g}(B_{\Theta})\) which packages higher-genus handle operators and exhibits a sharp dichotomy:
after inverting \(u\), higher genus becomes Morita-trivial; modulo \(u\), it degenerates to an exterior extension by \(H_{1}(\Sigma_{g})\).
\end{enumerate}
These statements are proved here and belong to what the monograph elsewhere calls the \emph{M-level}: strict algebraic models.

\medskip

\noindent\textbf{What remains interpretive or conjectural.}
Whenever we identify \(B\) with a bulk Koszul-dual algebra \(A^{!}\), or \(\Theta\) with a physical genus anomaly \(\Theta_{g}\), or \(\mathfrak{G}_{g}(B_{\Theta})\) with the full higher-genus line-side package of a twisted field theory, we cross from strict algebra into a physical programme.
Those identifications are supported by \cite{Zeng23,CDG20,DNP25,costello-gaiotto,CP2020}, but outside the strict models below they require caution.
The distinction is made explicit throughout.

\subsection{The first-principles correction}
% label removed: subsec:tholog-first-principles-correction

One might attempt to adjoin to \(B\) an odd \emph{supercentral} generator \(\eta\) with \(d\eta=\Theta\).
This does not produce the correct object.
Indeed, if \(\eta\) were odd and supercentral in characteristic \(0\), then
\[
\eta^{2}=-\eta^{2},
\]
hence \(\eta^{2}=0\), and the entire secondary theory would collapse before it begins.
The higher-genus refinement would then have no place to live.

The correct object is a \emph{skew polynomial extension}:
\(\eta\) graded-commutes with \(B\) but remains a free odd generator with respect to its own powers.
Equivalently, one adjoins an odd variable over \(B\) via the sign automorphism
\[
\sigma:B\longrightarrow B,\qquad \sigma(b)=(-1)^{|b|}b.
\]
This is the smallest repair preserving the universal property while leaving room for a nontrivial square
\[
u:=\eta^{2}.
\]

A second correction is equally important.
Because the extension is a skew/Ore extension, the correct exact triangle is \emph{one-sided} and module-theoretic, not the naive untwisted bimodule resolution.
The replacement is a two-term free resolution of each module, from which the transgression formulas for \(R\!\operatorname{Hom}\) and derived tensor products follow directly.

\begin{remark}[Why this construction is inevitable]
% label removed: rem:tholog-inevitability
The transgression algebra $B_\Theta$ is not one possible
resolution of the curvature obstruction among several: it is
the \emph{universal} one. Its universal property
(Proposition~\ref{prop:tholog-universal-property}) states that
every dg~algebra map $B_\Theta \to C$ is equivalent to a pair
$(f\colon B \to C,\; h \in C^1)$ with $dh = f(\Theta)$ and
$hf(b) = (-1)^{|b|}f(b)h$. In other words, $B_\Theta$
\emph{is} the moduli problem ``resolve $\Theta$ at chain
level'', incarnated as an algebra. The secondary anomaly
$u = \eta^2$ and the genus-Clifford completion
$\mathfrak G_g(B_\Theta)$ are then forced by the same
universality: $u$ is the inevitable by-product of
$\eta^2 \ne 0$, and the Clifford algebra is the universal
algebra in which $2g$ odd generators square to~$u$. The
entire tower (transgression, neutralisation, secondary
anomaly, genus completion, Morita triviality) is a single
deductive chain from the datum $(B, \Theta)$. No choices are
made; no alternatives exist at the same level of universality.
\end{remark}

\subsection{Conventions}
% label removed: subsec:tholog-conventions

Throughout this appendix:
\begin{itemize}
\item all complexes are cohomologically graded;
\item the ground field \(k\) has characteristic \(0\);
\item differentials have degree \(+1\);
\item if \(X\) is a graded object, \(|x|\) denotes the degree of a homogeneous element \(x\in X\);
\item for a dg algebra \(B\), \(Z(B)\) denotes its graded center;
\item if \(M\) is a left dg \(B\)-module, \(\underline{\operatorname{Hom}}_{B}(M,N)\) denotes the dg Hom complex;
\item derived functors are written \(R\!\operatorname{Hom}\) and \(\otimes^{L}\);
\item \(\Sigma_{g}\) denotes a closed connected oriented surface of genus \(g\).
\end{itemize}

\subsection{From the source picture to the strict algebraic problem}
% label removed: subsec:tholog-source-picture

Let us compress the surrounding story into one strict question.

\begin{quote}
Given a dg algebra \(B\) modelling bulk, defect, or line-operator data, and given a closed central degree-\(2\) anomaly class
\(\Theta\in Z^{2}(B)\cap Z(B)\),
what is the correct universal algebra governing those modules, bimodules, and tensor products on which the action of \(\Theta\) has been neutralized at chain level?
How does one incorporate genus in a way that is invisible away from the anomaly locus and maximally nontrivial on it?
\end{quote}

The answer will be:
\begin{enumerate}
\item first pass from \(B\) to the transgression algebra \(B_{\Theta}\);
\item then, if one wishes to remember genuine surface data, pass further to the genus-Clifford completion \(\mathfrak{G}_{g}(B_{\Theta})\).
\end{enumerate}

Everything below is designed so that these two steps are universal, exact, functorial, and compatible with the familiar monoidal structures appearing in line-operator theories.

\subsection{The transgression algebra}
% label removed: subsec:tholog-transgression-algebra

Let \(B\) be a dg algebra, and let
\[
\Theta\in Z^{2}(B)\cap Z(B)
\]
be a closed central degree-\(2\) element.

\begin{definition}[Transgression algebra]
% label removed: def:tholog-transgression-algebra
The \emph{transgression algebra} of \((B,\Theta)\) is the dg algebra
\[
B_{\Theta}
:=
\frac{B * k\langle \eta\rangle}
{\langle \eta b-(-1)^{|b|}b\eta\;\mid\; b\in B \text{ homogeneous}\rangle},
\qquad |\eta|=1,\qquad d\eta=\Theta,
\]
where \(B * k\langle\eta\rangle\) denotes the free product of graded algebras and the differential extends that of \(B\).
\end{definition}

Thus \(B_{\Theta}\) is the skew polynomial/Ore extension of \(B\) by the sign automorphism \(\sigma(b)=(-1)^{|b|}b\), equipped with a differential transgressing \(\Theta\).

\begin{remark}[Normal form]
% label removed: rem:tholog-normal-form
Every element of \(B_{\Theta}\) may be written uniquely as a finite sum
\[
\sum_{n\geq 0} b_{n}\eta^{n},
\qquad b_{n}\in B,
\]
with all powers of \(\eta\) placed on the right.
Equivalently, \(B_{\Theta}\) is free as both a left and a right \(B\)-module on the basis \(\{1,\eta,\eta^{2},\dots\}\).
\end{remark}

\begin{proposition}[Universal property; \ClaimStatusProvedHere]
% label removed: prop:tholog-universal-property
Let \(C\) be a dg algebra.
To give a dg algebra map
\[
\widetilde f:B_{\Theta}\longrightarrow C
\]
is equivalent to giving:
\begin{enumerate}
\item a dg algebra map \(f:B\to C\);
\item a degree-\(1\) element \(h\in C^{1}\)
\end{enumerate}
such that
\[
h\,f(b)=(-1)^{|b|}f(b)\,h
\qquad \text{for all homogeneous } b\in B,
\]
and
\[
dh=f(\Theta).
\]
Under this correspondence one has \(\widetilde f(\eta)=h\).
\end{proposition}

\begin{proof}
A dg algebra map \(\widetilde f\) is determined by its restriction \(f=\widetilde f|_{B}\) and by the image
\(h=\widetilde f(\eta)\).
Preservation of the defining relations for \(B_{\Theta}\) is exactly the relation
\[
h\,f(b)=(-1)^{|b|}f(b)\,h,
\]
and compatibility with the differential is exactly
\[
dh=d\widetilde f(\eta)=\widetilde f(d\eta)=\widetilde f(\Theta)=f(\Theta).
\]
Conversely, any pair \((f,h)\) satisfying those two conditions extends uniquely to a dg algebra map on the free product and descends through the quotient by the relation \(\eta b=(-1)^{|b|}b\eta\).
\end{proof}

\begin{corollary}[Modules as neutralizations; \ClaimStatusProvedHere]
% label removed: cor:tholog-modules-neutralizations
To give a left dg \(B_{\Theta}\)-module structure on a graded vector space \(M\) is equivalent to giving:
\begin{enumerate}
\item a left dg \(B\)-module structure on \(M\);
\item a degree-\(1\) \(B\)-linear endomorphism
\[
h_{M}\in \underline{\operatorname{End}}^{1}_{B}(M)
\]
such that
\[
d_{\underline{\operatorname{End}}}(h_{M})=\Theta_{M},
\]
where \(\Theta_{M}\) denotes the action of \(\Theta\) on \(M\).
\end{enumerate}
\end{corollary}

\begin{proof}
Apply Proposition~\ref{prop:tholog-universal-property} to
\[
C=\underline{\operatorname{End}}(M)
\]
and note that the \(B\)-centrality condition becomes precisely \(B\)-linearity of the degree-\(1\) endomorphism \(h_{M}\).
\end{proof}

\begin{definition}[Neutralizable module]
% label removed: def:tholog-neutralizable
A left dg \(B\)-module \(M\) will be called \(\Theta\)-\emph{neutralizable} if the action of \(\Theta\) on \(M\) is null-homotopic, i.e.\ if there exists
\[
h_{M}\in \underline{\operatorname{End}}^{1}_{B}(M)
\]
with
\[
d_{\underline{\operatorname{End}}}(h_{M})=\Theta_{M}.
\]
Any such \(h_{M}\) will be called a \emph{neutralization} of the \(\Theta\)-action on \(M\).
\end{definition}

\begin{corollary}[Obstruction class; \ClaimStatusProvedHere]
% label removed: cor:tholog-obstruction-class
A left dg \(B\)-module \(M\) is \(\Theta\)-neutralizable if and only if
\[
[\Theta_{M}]=0\in H^{2}\underline{\operatorname{End}}_{B}(M)
=\operatorname{Ext}^{2}_{B}(M,M).
\]
\end{corollary}

\begin{proof}
By definition, \(M\) is neutralizable if and only if \(\Theta_{M}\) is a coboundary in the dg algebra \(\underline{\operatorname{End}}_{B}(M)\).
\end{proof}

\begin{corollary}[Moduli of neutralizations; \ClaimStatusProvedHere]
% label removed: cor:tholog-moduli-neutralizations
If \(M\) is \(\Theta\)-neutralizable, then the set of homotopy classes of neutralizations of \(M\) is an affine space over
\[
H^{1}\underline{\operatorname{End}}_{B}(M)
=\operatorname{Ext}^{1}_{B}(M,M).
\]
\end{corollary}

\begin{proof}
If \(h_{M}\) and \(h_{M}'\) are two neutralizations, then
\[
d(h_{M}'-h_{M})=0.
\]
Thus the difference is a degree-\(1\) cocycle.
Changing a neutralization by a boundary does not change its homotopy class.
\end{proof}

\subsection{The corrected one-sided resolution}
% label removed: subsec:tholog-corrected-one-sided-resolution

The next theorem is the corrected first-principles replacement for the naive untwisted bimodule resolution.
It is the technical engine behind the exact formulas for derived morphisms and relative tensor products.

\begin{theorem}[Free two-term resolution of a neutralized module; \ClaimStatusProvedHere]
% label removed: thm:tholog-free-two-term-resolution
Let \(M\) be a left dg \(B_{\Theta}\)-module, and let \(h_{M}\) be the corresponding neutralization of the underlying left dg \(B\)-module.
Define a degree-\(1\) map of graded left \(B_{\Theta}\)-modules
\[
D_{M}:B_{\Theta}\otimes_{B} M\longrightarrow B_{\Theta}\otimes_{B} M
\]
by
\[
D_{M}(s\otimes m)=s\eta\otimes m-(-1)^{|s|}s\otimes h_{M}(m).
\]
Then \(D_{M}\) is a chain map of degree \(1\), and there is a short exact sequence of dg \(B_{\Theta}\)-modules
\[
0\longrightarrow \big(B_{\Theta}\otimes_{B}M\big)[-1]
\xrightarrow{\;\widetilde D_{M}\;}
B_{\Theta}\otimes_{B}M
\xrightarrow{\;\mu_{M}\;}
M
\longrightarrow 0,
\]
where \(\widetilde D_{M}\) is the corresponding degree-\(0\) map from the shifted domain and \(\mu_{M}(s\otimes m)=s\cdot m\).
\end{theorem}

\begin{proof}
We first check that \(D_{M}\) is well defined over \(B\).
For homogeneous \(b\in B\),
\[
D_{M}(sb\otimes m)
=
sb\eta\otimes m-(-1)^{|s|+|b|}sb\otimes h_{M}(m).
\]
Since \(b\eta=(-1)^{|b|}\eta b\) in \(B_{\Theta}\),
\[
sb\eta\otimes m=(-1)^{|b|}s\eta b\otimes m=s\eta\otimes bm,
\]
and because \(h_{M}\) is \(B\)-linear,
\[
sb\otimes h_{M}(m)=s\otimes bh_{M}(m)=s\otimes h_{M}(bm).
\]
Hence
\[
D_{M}(sb\otimes m)=D_{M}(s\otimes bm).
\]

We next show that \(D_{M}\) is a chain map.
Write \(d\) for the differential on \(B_{\Theta}\otimes_{B}M\).
A direct computation gives
\begin{align*}
d\,D_{M}(s\otimes m)
&=
d(s\eta)\otimes m+(-1)^{|s|+1}s\eta\otimes dm
\\
&\qquad
-(-1)^{|s|}d(s)\otimes h_{M}(m)-s\otimes d(h_{M}m),
\end{align*}
while
\begin{align*}
D_{M}\,d(s\otimes m)
&=
D_{M}\big(d(s)\otimes m+(-1)^{|s|}s\otimes dm\big)
\\
&=
d(s)\eta\otimes m-(-1)^{|s|+1}d(s)\otimes h_{M}(m)
\\
&\qquad
+(-1)^{|s|}s\eta\otimes dm-s\otimes h_{M}(dm).
\end{align*}
Subtracting, using \(d(s\eta)=d(s)\eta+(-1)^{|s|}s\Theta\) and
\[
d(h_{M}m)=d_{\underline{\operatorname{End}}}(h_{M})(m)-h_{M}(dm)=\Theta_{M}(m)-h_{M}(dm),
\]
we obtain
\[
dD_{M}(s\otimes m)+D_{M}d(s\otimes m)=
(-1)^{|s|}s\Theta\otimes m-s\otimes \Theta_{M}(m)=0.
\]
So \(D_{M}\) is a chain map.

Surjectivity of \(\mu_{M}\) is obvious.
The image of \(D_{M}\) lies in \(\ker(\mu_{M})\) because
\[
\mu_{M}\big(D_{M}(s\otimes m)\big)
=s\eta\cdot m-(-1)^{|s|}s\cdot h_{M}(m)=0.
\]
Conversely, \(B_{\Theta}\otimes_{B}M\) is the free left \(B_{\Theta}\)-module on the underlying \(B\)-module \(M\), and \(M\) is obtained from it by imposing the relations
\[
\eta\otimes m=1\otimes h_{M}(m).
\]
Hence \(\ker(\mu_{M})\) is generated by the image of \(D_{M}\).

It remains to prove injectivity of \(\widetilde D_{M}\), equivalently of \(D_{M}\) as a degree-\(1\) map on the underlying graded modules.
Write an element in normal form:
\[
x=\sum_{i=0}^{N}\eta^{i}\otimes m_{i}.
\]
Then
\[
D_{M}(x)=\sum_{i=0}^{N}\eta^{i+1}\otimes m_{i}-\sum_{i=0}^{N}(-1)^{i}\eta^{i}\otimes h_{M}(m_{i}).
\]
The coefficient of \(\eta^{N+1}\) is \(m_{N}\), so if \(D_{M}(x)=0\), then \(m_{N}=0\).
Descending induction on \(N\) gives \(m_{i}=0\) for all \(i\).
Thus \(D_{M}\) is injective.
\end{proof}

\begin{remark}[Why one-sided is the right level]
% label removed: rem:tholog-why-one-sided
Theorem~\ref{thm:tholog-free-two-term-resolution} is the precise place where the Ore nature of \(B_{\Theta}\) matters.
One should think of \(B_{\Theta}\) not as \(B\) tensored with an exterior algebra, but as a skew polynomial extension.
The resulting exact triangle is naturally attached to a \emph{module}, not first to a bimodule.
This is the correct algebraic setting for the transgression formulas below.
\end{remark}

\subsection{Transgression of derived morphisms}
% label removed: subsec:tholog-derived-morphisms

\begin{theorem}[Transgression exact triangle for derived Hom; \ClaimStatusProvedHere]
% label removed: thm:tholog-derived-hom
Let \(M,N\) be left dg \(B_{\Theta}\)-modules, with neutralizations \(h_{M}\) and \(h_{N}\).
Define the degree-\(1\) endomorphism
\[
\delta_{M,N}:\underline{\operatorname{Hom}}_{B}(M,N)\longrightarrow
\underline{\operatorname{Hom}}_{B}(M,N)
\]
by
\[
\delta_{M,N}(f)=h_{N}\circ f-(-1)^{|f|}f\circ h_{M}.
\]
Then there is a functorial exact triangle
\[
R\!\operatorname{Hom}_{B_{\Theta}}(M,N)
\longrightarrow
R\!\operatorname{Hom}_{B}(M,N)
\xrightarrow{\;\delta_{M,N}\;}
R\!\operatorname{Hom}_{B}(M,N)[1]
\longrightarrow .
\]
Equivalently,
\[
R\!\operatorname{Hom}_{B_{\Theta}}(M,N)
\simeq
\operatorname{Fib}\!\left(
R\!\operatorname{Hom}_{B}(M,N)
\xrightarrow{\;\delta_{M,N}\;}
R\!\operatorname{Hom}_{B}(M,N)[1]
\right).
\]
\end{theorem}

\begin{proof}
Apply \(\underline{\operatorname{Hom}}_{B_{\Theta}}(-,N)\) to the short exact sequence of Theorem~\ref{thm:tholog-free-two-term-resolution}.
Because \(B_{\Theta}\otimes_{B}M\) is a free left \(B_{\Theta}\)-module on the underlying \(B\)-module \(M\), the standard adjunction gives
\[
\underline{\operatorname{Hom}}_{B_{\Theta}}(B_{\Theta}\otimes_{B}M,N)
\cong
\underline{\operatorname{Hom}}_{B}(M,N).
\]
Under this identification, precomposition by \(D_{M}\) becomes exactly
\[
f\longmapsto h_{N}f-(-1)^{|f|}fh_{M}.
\]
Passing to derived categories yields the exact triangle.
\end{proof}

\begin{corollary}[Long exact sequence in Ext; \ClaimStatusProvedHere]
% label removed: cor:tholog-long-exact-ext
For any left dg \(B_{\Theta}\)-modules \(M,N\), there is a natural long exact sequence
\[
\cdots\to
\operatorname{Ext}^{n}_{B_{\Theta}}(M,N)
\to
\operatorname{Ext}^{n}_{B}(M,N)
\xrightarrow{\;\delta_{*}\;}
\operatorname{Ext}^{n+1}_{B}(M,N)
\to
\operatorname{Ext}^{n+1}_{B_{\Theta}}(M,N)
\to\cdots.
\]
\end{corollary}

\begin{remark}[The philosophy]
% label removed: rem:tholog-hom-philosophy
The exact triangle says that the anomaly-completed morphism complex is not mysterious.
It is computed from the underlying \(B\)-morphism complex by imposing one explicit degree-\(1\) transgression operator \(\delta_{M,N}\).
In particular, anomaly-completed Ext groups are controlled by ordinary Ext groups plus one universal connecting morphism.
\end{remark}

\subsection{Transgression of derived relative tensor products}
% label removed: subsec:tholog-derived-tensors

For a right \(B_{\Theta}\)-module, the neutralization operator is not strictly right \(B\)-linear; it is twisted by the sign automorphism.
We therefore record the correct right-handed notion.

\begin{definition}[Right neutralization]
% label removed: def:tholog-right-neutralization
A right dg \(B_{\Theta}\)-module \(L\) is equivalently a right dg \(B\)-module together with a degree-\(1\) endomorphism
\[
h_{L}:L\longrightarrow L
\]
such that
\[
h_{L}(lb)=(-1)^{|b|}h_{L}(l)b
\qquad (b\in B \text{ homogeneous}),
\]
and
\[
d(h_{L})=\Theta_{L},
\]
where \(\Theta_{L}\) denotes right multiplication by \(\Theta\).
We call \(h_{L}\) a \emph{right neutralization}.
\end{definition}

\begin{theorem}[Transgression exact triangle for derived tensor products; \ClaimStatusProvedHere]
% label removed: thm:tholog-derived-tensor
Let \(L\) be a right dg \(B_{\Theta}\)-module and \(R\) a left dg \(B_{\Theta}\)-module, with right and left neutralizations \(h_{L}\) and \(h_{R}\).
Then there is a functorial exact triangle
\[
L\otimes^{L}_{B}R[-1]
\xrightarrow{\;\theta_{L,R}\;}
L\otimes^{L}_{B}R
\longrightarrow
L\otimes^{L}_{B_{\Theta}}R
\longrightarrow ,
\]
where on homogeneous tensors
\[
\theta_{L,R}(l\otimes r)
=
(-1)^{|l|}h_{L}(l)\otimes r-l\otimes h_{R}(r).
\]
Equivalently,
\[
L\otimes^{L}_{B_{\Theta}}R
\simeq
\operatorname{Cone}\!\left(
L\otimes^{L}_{B}R[-1]
\xrightarrow{\;\theta_{L,R}\;}
L\otimes^{L}_{B}R
\right).
\]
\end{theorem}

\begin{proof}
Resolve \(R\) by the free two-term resolution from Theorem~\ref{thm:tholog-free-two-term-resolution}, then tensor on the left over \(B_{\Theta}\) with \(L\).
The resulting two-term complex is
\[
L\otimes_{B}R[-1]\xrightarrow{\theta_{L,R}}L\otimes_{B}R,
\]
and a direct inspection shows that the induced map is exactly
\[
l\otimes r\longmapsto (-1)^{|l|}h_{L}(l)\otimes r-l\otimes h_{R}(r).
\]
The sign in the first term is forced by the \(\sigma\)-semilinearity of the right neutralization and ensures that the map is \(B\)-balanced.
Passing to derived tensor products yields the triangle.
\end{proof}

\begin{remark}[Interval compactification]
% label removed: rem:tholog-interval-compactification
Theorem~\ref{thm:tholog-derived-tensor} is the strict algebraic shadow of the interval/strip compactification formula familiar from boundary topological field theory.
At the level of strict dg models, the first anomaly correction to an interval compactification is always a universal two-term cone.
This is exactly the form one would want when comparing boundary conditions through a bulk or line algebra endowed with a nontrivial central genus class.
\end{remark}

\subsection{Gauge invariance of the transgression operator}
% label removed: subsec:tholog-gauge-invariance

Neutralizations are never unique.
The next theorem says that the transgression operator on morphism complexes depends only on the \(B_{\Theta}\)-module structures, not on the chosen representatives.

\begin{theorem}[Gauge invariance of transgression; \ClaimStatusProvedHere]
% label removed: thm:tholog-gauge-invariance
Let \(M,N\) be fixed left dg \(B\)-modules, and suppose
\[
h_{M}'=h_{M}+d(u_{M}),
\qquad
h_{N}'=h_{N}+d(u_{N}),
\]
for degree-\(0\) \(B\)-linear endomorphisms \(u_{M},u_{N}\).
Then the corresponding operators
\[
\delta_{M,N},\delta_{M,N}' :
\underline{\operatorname{Hom}}_{B}(M,N)\to
\underline{\operatorname{Hom}}_{B}(M,N)[1]
\]
are chain-homotopic.
\end{theorem}

\begin{proof}
Define a degree-\(0\) endomorphism
\[
H(f)=u_{N}\circ f-f\circ u_{M}.
\]
A direct computation in the dg algebra \(\underline{\operatorname{Hom}}_{B}(M,N)\) gives
\[
\delta_{M,N}'-\delta_{M,N}=dH+Hd.
\]
Hence the two operators are chain-homotopic.
\end{proof}

\subsection{Monoidal lift for primitive anomalies}
% label removed: subsec:tholog-monoidal-lift

The next theorem explains when the anomaly-completed theory inherits a monoidal structure.

\begin{definition}[Primitive anomaly]
% label removed: def:tholog-primitive-anomaly
Assume that \(B\) is a dg bialgebra with coproduct \(\Delta\) and counit \(\varepsilon\).
A closed central degree-\(2\) class
\[
\Theta\in Z^{2}(B)\cap Z(B)
\]
will be called \emph{primitive} if
\[
\Delta(\Theta)=\Theta\otimes 1+1\otimes \Theta,
\qquad
\varepsilon(\Theta)=0.
\]
\end{definition}

\begin{theorem}[Primitive lift theorem; \ClaimStatusProvedHere]
% label removed: thm:tholog-primitive-lift
Let \(B\) be a dg bialgebra and let \(\Theta\in Z^{2}(B)\cap Z(B)\) be primitive.
Then the transgression algebra \(B_{\Theta}\) admits a unique dg bialgebra structure extending that of \(B\) and satisfying
\[
\Delta(\eta)=\eta\otimes 1+1\otimes \eta,
\qquad
\varepsilon(\eta)=0.
\]
\end{theorem}

\begin{proof}
Because \(B_{\Theta}\) is generated as a graded algebra by \(B\) and \(\eta\), uniqueness is immediate.

Define \(\Delta\) and \(\varepsilon\) on generators by the given formulas.
It remains to check that the defining relation
\[
\eta b=(-1)^{|b|}b\eta
\]
is preserved.
Write
\[
\Delta(b)=\sum b_{(1)}\otimes b_{(2)}.
\]
Then in the graded tensor product algebra,
\begin{align*}
\Delta(\eta)\Delta(b)
&=
(\eta\otimes 1+1\otimes \eta)\sum b_{(1)}\otimes b_{(2)}
\\
&=
\sum \eta b_{(1)}\otimes b_{(2)}+\sum (-1)^{|b_{(1)}|}b_{(1)}\otimes \eta b_{(2)}
\\
&=
\sum (-1)^{|b_{(1)}|}b_{(1)}\eta\otimes b_{(2)}
+\sum (-1)^{|b|}b_{(1)}\otimes b_{(2)}\eta
\\
&=
(-1)^{|b|}\Delta(b)\Delta(\eta).
\end{align*}
So the relation is respected.

Compatibility with the differential on \(\eta\) is the primitive condition:
\[
d\,\Delta(\eta)
=
\Theta\otimes 1+1\otimes \Theta
=
\Delta(\Theta)
=
\Delta(d\eta).
\]
All other identities are inherited from \(B\).
\end{proof}

\begin{corollary}[Fusion of neutralizations; \ClaimStatusProvedHere]
% label removed: cor:tholog-fusion-neutralizations
Assume the setting of Theorem~\ref{thm:tholog-primitive-lift}.
If \(M,N\) are left dg \(B_{\Theta}\)-modules with neutralizations \(h_{M}\) and \(h_{N}\), then \(M\otimes N\) becomes a left dg \(B_{\Theta}\)-module with neutralization
\[
h_{M\otimes N}=h_{M}\otimes 1+1\otimes h_{N},
\]
where \(1\otimes h_{N}\) is understood in the graded sense, so that on homogeneous tensors
\[
h_{M\otimes N}(m\otimes n)=h_{M}(m)\otimes n+(-1)^{|m|}m\otimes h_{N}(n).
\]
\end{corollary}

\begin{proof}
This is the action induced by \(\Delta(\eta)=\eta\otimes 1+1\otimes \eta\).
Since \(\eta\) is primitive, the differential condition follows from \(d(h_{M})=\Theta_{M}\) and \(d(h_{N})=\Theta_{N}\).
\end{proof}

\subsection{The secondary anomaly}
% label removed: subsec:tholog-secondary-anomaly

The new ingredient which is invisible in naive supercentral models and inevitable in the corrected skew model is the central square of the transgression generator.

\begin{proposition}[The secondary central class; \ClaimStatusProvedHere]
% label removed: prop:tholog-secondary-central-class
Let
\[
u:=\eta^{2}\in B_{\Theta}^{2}.
\]
Then \(u\) is a closed central degree-\(2\) element of \(B_{\Theta}\).
\end{proposition}

\begin{proof}
Since \(d\eta=\Theta\) and \(\Theta\) has even degree,
\[
d(u)=d(\eta^{2})=(d\eta)\eta-\eta(d\eta)=\Theta\eta-\eta\Theta=0,
\]
because \(\eta\) graded-commutes with \(B\), hence with \(\Theta\).

To check centrality, let \(b\in B\) be homogeneous.
Then
\[
u\,b=\eta(\eta b)=(-1)^{|b|}\eta(b\eta)=(-1)^{2|b|}b\eta^{2}=bu.
\]
Clearly \(u\) commutes with \(\eta\) by associativity, hence with all of \(B_{\Theta}\).
\end{proof}

\begin{definition}[Secondary anomaly of a lift]
% label removed: def:tholog-secondary-anomaly
Let \(M\) be a \(\Theta\)-neutralizable left dg \(B\)-module and let \(h\) be a chosen neutralization.
The \emph{secondary anomaly} of the lift \((M,h)\) is the cohomology class
\[
\sigma(M,h):=[h^{2}]
\in
H^{2}\underline{\operatorname{End}}_{B}(M)
=
\operatorname{Ext}^{2}_{B}(M,M).
\]
Equivalently, \(\sigma(M,h)\) is the action of the central element \(u=\eta^{2}\) on \(M\).
\end{definition}

\begin{remark}[Why this is genuinely new]
% label removed: rem:tholog-secondary-why-new
The secondary anomaly is invisible if one forces \(\eta\) to be odd supercentral.
In the corrected skew extension, however, \(u=\eta^{2}\) is honest, closed, central, and nontrivial.
It is exactly the place where higher genus will land.
\end{remark}

\subsection{Quadratic geometry of the space of neutralizations}
% label removed: subsec:tholog-quadratic-geometry

The moduli of neutralizations is not merely affine.
It carries a canonical quadratic refinement valued in degree-\(2\) self-Ext.

\begin{theorem}[Quadratic refinement on the space of neutralizations; \ClaimStatusProvedHere]
% label removed: thm:tholog-quadratic-refinement
Let \(M\) be a \(\Theta\)-neutralizable left dg \(B\)-module, and choose one neutralization \(h\).
Let
\[
\mathcal{L}_{\Theta}(M)
:=
\big\{
h'\in \underline{\operatorname{End}}^{1}_{B}(M)\;\big|\;
d(h')=\Theta_{M}
\big\}.
\]
Then:
\begin{enumerate}
\item \(\mathcal{L}_{\Theta}(M)\) is an affine space over the degree-\(1\) cocycles
\[
Z^{1}\underline{\operatorname{End}}_{B}(M).
\]
\item The assignment
\[
q_{h}(z):=[(h+z)^{2}-h^{2}]
\]
descends to a well-defined map
\[
q_{h}:H^{1}\underline{\operatorname{End}}_{B}(M)\longrightarrow
H^{2}\underline{\operatorname{End}}_{B}(M),
\]
and admits the explicit formula
\[
q_{h}([z])=[hz+zh+z^{2}].
\]
\item Its polarization
\[
b([z],[w])
=
q_{h}([z]+[w])-q_{h}([z])-q_{h}([w])
\]
is independent of the chosen base point \(h\) and is given by
\[
b([z],[w])=[zw+wz].
\]
Thus \(H^{1}\underline{\operatorname{End}}_{B}(M)\) carries a canonical symmetric bilinear form with values in \(H^{2}\underline{\operatorname{End}}_{B}(M)\), and each neutralization \(h\) provides a quadratic refinement of that form.
\end{enumerate}
\end{theorem}

\begin{proof}
If \(h',h''\in \mathcal{L}_{\Theta}(M)\), then
\[
d(h''-h')=0,
\]
so \(\mathcal{L}_{\Theta}(M)\) is an affine space over \(Z^{1}\underline{\operatorname{End}}_{B}(M)\).
This proves (1).

For (2), if \(z\in Z^{1}\underline{\operatorname{End}}_{B}(M)\), then \(h+z\) is again a neutralization, so \((h+z)^{2}-h^{2}=hz+zh+z^{2}\) is a degree-\(2\) cocycle.
We must show that its cohomology class depends only on \([z]\).
Replace \(z\) by \(z+d(a)\) with \(a\in \underline{\operatorname{End}}^{0}_{B}(M)\).
Set \(x=h+z\), so \(dx=\Theta_{M}\).
Then
\[
(x+d(a))^{2}-x^{2}=x\,d(a)+d(a)\,x+d(a)^{2}.
\]
Since \(d^{2}(a)=[\Theta_{M},a]=0\) (because \(\Theta_{M}\) comes from the action of the central element \(\Theta\in B\)), one computes
\[
x\,d(a)+d(a)\,x+d(a)^{2}
=
d\big(-xa+ax+a\,d(a)\big).
\]
Thus the cohomology class is unchanged.
This proves (2).

Finally, polarization gives
\[
q_{h}(z+w)-q_{h}(z)-q_{h}(w)=zw+wz,
\]
which is independent of \(h\).
This proves (3).
\end{proof}

\subsection{Curvature identity on morphism complexes}
% label removed: subsec:tholog-curvature-identity

The transgression operator \(\delta_{M,N}\) is rarely square-zero.
Its failure to square to zero is exactly measured by the secondary anomaly.

\begin{theorem}[Curvature identity; \ClaimStatusProvedHere]
% label removed: thm:tholog-curvature-identity
Let \(M,N\) be left dg \(B_{\Theta}\)-modules.
Then the operator
\[
\delta_{M,N}(f)=h_{N}f-(-1)^{|f|}fh_{M}
\]
from Theorem~\ref{thm:tholog-derived-hom} is a degree-\(1\) chain map satisfying
\[
\delta_{M,N}^{2}(f)=h_{N}^{2}f-fh_{M}^{2}.
\]
Equivalently,
\[
\delta_{M,N}^{2}
=
u_{N}\circ(-)-(-)\circ u_{M},
\]
where \(u_{M}=h_{M}^{2}\) and \(u_{N}=h_{N}^{2}\).
Passing to cohomology,
\[
[\delta_{M,N}]^{2}
=
\sigma(N,h_{N})\circ(-)-(-)\circ \sigma(M,h_{M})
\]
on \(\operatorname{Ext}^{\bullet}_{B}(M,N)\).
\end{theorem}

\begin{proof}
A direct computation shows that \(\delta_{M,N}\) is a chain map because
\[
d(h_{N})=\Theta_{N},\qquad d(h_{M})=\Theta_{M},
\]
and the \(\Theta\)-terms cancel by centrality.
Squaring \(\delta_{M,N}\), we obtain
\begin{align*}
\delta_{M,N}^{2}(f)
&=
h_{N}\big(h_{N}f-(-1)^{|f|}fh_{M}\big)
-
(-1)^{|f|+1}\big(h_{N}f-(-1)^{|f|}fh_{M}\big)h_{M}
\\
&=
h_{N}^{2}f-fh_{M}^{2},
\end{align*}
because the middle terms cancel.
The cohomological statement follows immediately.
\end{proof}

\begin{corollary}[Orthogonality of sectors with distinct secondary anomaly; \ClaimStatusProvedHere]
% label removed: cor:tholog-orthogonality
Let \(M,N\) be left dg \(B_{\Theta}\)-modules.
If the degree-\(2\) endomorphism
\[
u_{N}\circ(-)-(-)\circ u_{M}
\]
acts by a homotopy-invertible endomorphism on \(R\!\operatorname{Hom}_{B}(M,N)\), then
\[
R\!\operatorname{Hom}_{B_{\Theta}}(M,N)\simeq 0.
\]
\end{corollary}

\begin{proof}
By Theorem~\ref{thm:tholog-curvature-identity}, the transgression operator \(\delta_{M,N}\) has square equal to the displayed endomorphism.
If that square is homotopy invertible, then \(\delta_{M,N}\) itself is homotopy invertible, with homotopy inverse
\[
\big(u_{N}\circ(-)-(-)\circ u_{M}\big)^{-1}\delta_{M,N}.
\]
But the fiber of a homotopy equivalence is acyclic, so Theorem~\ref{thm:tholog-derived-hom} gives the claim.
\end{proof}

\begin{corollary}[Equal-secondary sector; \ClaimStatusProvedHere]
% label removed: cor:tholog-equal-secondary-sector
If \(u_{M}=u_{N}\) as chain maps, then
\[
\delta_{M,N}^{2}=0.
\]
Hence \(\underline{\operatorname{Hom}}_{B}(M,N)\) becomes a mixed complex with differentials \(d\) and \(\delta_{M,N}\), and one has a short exact sequence
\[
0\longrightarrow
\operatorname{coker}\!\big(H^{n-1}\delta_{M,N}\big)
\longrightarrow
\operatorname{Ext}^{n}_{B_{\Theta}}(M,N)
\longrightarrow
\ker\!\big(H^{n}\delta_{M,N}\big)
\longrightarrow 0.
\]
\end{corollary}

\subsection{Cohomology of the transgression algebra as a hypersurface}
% label removed: subsec:tholog-hypersurface

The transgression algebra admits a very concrete cohomological description.

Let
\[
C:=B[u],\qquad |u|=2,\qquad du=0.
\]
As a graded left \(C\)-module, every element of \(B_{\Theta}\) has a unique decomposition
\[
c_{0}+c_{1}\eta,
\qquad c_{0},c_{1}\in C.
\]

\begin{theorem}[Hypersurface long exact sequence; \ClaimStatusProvedHere]
% label removed: thm:tholog-hypersurface
There is a short exact sequence of complexes of left \(C\)-modules
\[
0\longrightarrow C
\longrightarrow B_{\Theta}
\longrightarrow C[-1]
\longrightarrow 0,
\]
whose connecting morphism on cohomology is multiplication by \([\Theta]\) (up to the standard cohomological sign convention).
Consequently, there is a long exact sequence
\[
\cdots\to
H^{n}(C)\to H^{n}(B_{\Theta})\to H^{n-1}(C)
\xrightarrow{\cdot[\Theta]}
H^{n+1}(C)\to\cdots.
\]
If multiplication by \([\Theta]\) is injective on \(H^{\bullet}(B)[u]\), then
\[
H^{\bullet}(B_{\Theta})\cong H^{\bullet}(B)[u]/([\Theta]).
\]
\end{theorem}

\begin{proof}
The inclusion \(C\hookrightarrow B_{\Theta}\) is the inclusion of the even part \(\sum b_{n}u^{n}\), and the quotient identifies with \(C[-1]\) via
\[
c\longmapsto c\eta.
\]
So there is a short exact sequence of graded \(C\)-modules.
The differential on \(c\eta\) is
\[
d(c\eta)=d(c)\eta+(-1)^{|c|}c\Theta.
\]
Hence the connecting morphism sends the class of \(c\) to the class of \(c\Theta\), up to the standard sign from the shift.
This yields the long exact sequence.
If multiplication by \([\Theta]\) is injective, exactness identifies \(H^{\bullet}(B_{\Theta})\) with the quotient of \(H^{\bullet}(C)=H^{\bullet}(B)[u]\) by the principal ideal generated by \([\Theta]\).
\end{proof}

\begin{remark}[The hypersurface viewpoint]
% label removed: rem:tholog-hypersurface-viewpoint
Theorem~\ref{thm:tholog-hypersurface} says that anomaly completion replaces \(B\) by a one-parameter degree-\(2\) hypersurface thickening.
The variable \(u\) is not an auxiliary decoration.
It is the ambient coordinate in which \([\Theta]\) cuts out the completed cohomology ring.
\end{remark}

\subsection{Genus-Clifford completion}
% label removed: subsec:tholog-genus-clifford

Adding genuine surface data.
A genus-\(g\) surface carries \(2g\) basic one-cycles, and in the anomaly-completed setting their algebraic shadow is an odd Clifford algebra whose anticommutator is controlled by the secondary class \(u\).

Let
\[
H:=H_{1}(\Sigma_{g};k).
\]
Its intersection pairing is symplectic.
After parity shift, \(V:=\Pi H\) carries the corresponding symmetric bilinear form.
Choose a symplectic basis
\[
a_{1},b_{1},\dots,a_{g},b_{g},
\]
and write
\[
\alpha_{i}:=\Pi a_{i},\qquad \beta_{i}:=\Pi b_{i}.
\]

\begin{definition}[Genus-Clifford completion]
% label removed: def:tholog-genus-clifford
The \emph{genus-\(g\) Clifford completion} of \(B_{\Theta}\) is the dg algebra
\[
\mathfrak{G}_{g}(B_{\Theta})
:=
\frac{
B_{\Theta}\langle \alpha_{1},\beta_{1},\dots,\alpha_{g},\beta_{g}\rangle
}{
\left\langle
\alpha_{i}\alpha_{j}+\alpha_{j}\alpha_{i},\;
\beta_{i}\beta_{j}+\beta_{j}\beta_{i},\;
\alpha_{i}\beta_{j}+\beta_{j}\alpha_{i}-\delta_{ij}u
\right\rangle
},
\]
with
\[
|\alpha_{i}|=|\beta_{i}|=1,\qquad d\alpha_{i}=d\beta_{i}=0.
\]
\end{definition}

\begin{remark}[Two regimes]
% label removed: rem:tholog-two-regimes
Because the odd generators satisfy a Clifford relation with coefficient \(u\), the genus theory naturally splits into two regimes:
\begin{enumerate}
\item the \emph{curved} regime \(u\neq 0\), in which the Clifford algebra becomes Morita-trivial after localization;
\item the \emph{strict} regime \(u=0\), in which the handle algebra degenerates to an exterior algebra on \(H_{1}(\Sigma_{g})\).
\end{enumerate}
This dichotomy is one of the main structural consequences of the whole construction.
\end{remark}

\begin{theorem}[Universal property of genus completion; \ClaimStatusProvedHere]
% label removed: thm:tholog-universal-property-genus
To give a left dg \(\mathfrak{G}_{g}(B_{\Theta})\)-module is equivalent to giving:
\begin{enumerate}
\item a left dg \(B_{\Theta}\)-module \(M\);
\item closed degree-\(1\) \(B_{\Theta}\)-linear endomorphisms
\[
A_{1},B_{1},\dots,A_{g},B_{g}\in \underline{\operatorname{End}}^{1}_{B_{\Theta}}(M)
\]
such that
\[
A_{i}A_{j}+A_{j}A_{i}=0,\qquad
B_{i}B_{j}+B_{j}B_{i}=0,\qquad
A_{i}B_{j}+B_{j}A_{i}=\delta_{ij}u_{M}.
\]
\end{enumerate}
\end{theorem}

\begin{proof}
This is immediate from generators and relations.
A module over \(\mathfrak{G}_{g}(B_{\Theta})\) is determined by the images of the generators
\(\alpha_{i},\beta_{i}\).
Preservation of the differential gives closedness of the corresponding endomorphisms, and preservation of the defining relations gives exactly the three displayed identities.
Conversely, any such family extends uniquely to a dg representation.
\end{proof}

\subsection{Morita triviality away from the secondary-anomaly cone}
% label removed: subsec:tholog-morita-triviality

Let
\[
R:=B_{\Theta}[u^{-1}]
\]
be the localization of \(B_{\Theta}\) at the central element \(u\).

\begin{theorem}[One handle becomes a matrix factor after inverting \(u\); \ClaimStatusProvedHere]
% label removed: thm:tholog-one-handle-matrix
There is an isomorphism of dg algebras
\[
\mathfrak{G}_{1}(R)\cong \underline{\operatorname{End}}_{R}
\big(R\oplus R\varepsilon\big),
\]
where \(|\varepsilon|=-1\).
\end{theorem}

\begin{proof}
Let
\[
S_{1}:=R\oplus R\varepsilon=R\otimes \Lambda^{\bullet}(k\varepsilon[-1]).
\]
Define degree-\(1\) endomorphisms of \(S_{1}\) by
\[
\rho(\alpha)=u\,(\varepsilon\wedge -),\qquad
\rho(\beta)=\iota_{\varepsilon}.
\]
Because
\[
(\varepsilon\wedge -)^{2}=0,\qquad \iota_{\varepsilon}^{2}=0,\qquad
(\varepsilon\wedge -)\iota_{\varepsilon}+\iota_{\varepsilon}(\varepsilon\wedge -)=\operatorname{id},
\]
we get
\[
\rho(\alpha)^{2}=0,\qquad
\rho(\beta)^{2}=0,\qquad
\rho(\alpha)\rho(\beta)+\rho(\beta)\rho(\alpha)=u\,\operatorname{id}.
\]
Hence \(\rho\) extends to a dg algebra map
\[
\mathfrak{G}_{1}(R)\longrightarrow \underline{\operatorname{End}}_{R}(S_{1}).
\]

As a left \(R\)-module, \(\mathfrak{G}_{1}(R)\) is free of rank \(4\) with basis
\[
1,\alpha,\beta,\alpha\beta.
\]
Likewise \(\underline{\operatorname{End}}_{R}(S_{1})\) is free of rank \(4\), with basis
\[
\operatorname{id},\quad \varepsilon\wedge -,\quad \iota_{\varepsilon},\quad
(\varepsilon\wedge -)\iota_{\varepsilon}.
\]
The image of \(\rho\) contains all four basis elements, because \(u\) is invertible and
\[
\varepsilon\wedge - = u^{-1}\rho(\alpha),\qquad
\iota_{\varepsilon}=\rho(\beta),\qquad
(\varepsilon\wedge -)\iota_{\varepsilon}=u^{-1}\rho(\alpha\beta).
\]
Thus \(\rho\) is an isomorphism.
\end{proof}

\begin{theorem}[Morita triviality of higher genus after inverting \(u\); \ClaimStatusProvedHere]
% label removed: thm:tholog-morita-triviality
For every \(g\geq 0\),
\[
\mathfrak{G}_{g}(B_{\Theta})[u^{-1}]
\cong
\underline{\operatorname{End}}_{R}
\Big(R\otimes \Lambda^{\bullet}(k^{g}[-1])\Big),
\qquad
R=B_{\Theta}[u^{-1}].
\]
In particular,
\[
D\!\big(\mathfrak{G}_{g}(B_{\Theta})[u^{-1}]\text{-mod}\big)
\simeq
D\!\big(B_{\Theta}[u^{-1}]\text{-mod}\big).
\]
\end{theorem}

\begin{proof}
By construction, \(\mathfrak{G}_{g}(R)\) is the graded tensor product of \(g\) copies of \(\mathfrak{G}_{1}(R)\).
Likewise
\[
R\otimes \Lambda^{\bullet}(k^{g}[-1])
\cong
S_{1}^{\otimes g}.
\]
Taking graded tensor products of the isomorphism from Theorem~\ref{thm:tholog-one-handle-matrix} yields
\[
\mathfrak{G}_{g}(R)\cong
\underline{\operatorname{End}}_{R}(S_{1}^{\otimes g})
\cong
\underline{\operatorname{End}}_{R}
\Big(R\otimes \Lambda^{\bullet}(k^{g}[-1])\Big).
\]
Derived Morita equivalence is immediate.
\end{proof}

\begin{corollary}[Localization principle; \ClaimStatusProvedHere]
% label removed: cor:tholog-localization-principle
All genuinely new genus-\(g\) information in \(\mathfrak{G}_{g}(B_{\Theta})\) is supported on the secondary-anomaly cone
\[
V(u)\subset \operatorname{Spec}H^{\bullet}(B_{\Theta}).
\]
Away from \(u=0\), the genus completion is derived Morita-trivial.
\end{corollary}

\subsection{Degeneration on the strict locus}
% label removed: subsec:tholog-degeneration-strict-locus

\begin{theorem}[Exterior degeneration modulo \(u\); \ClaimStatusProvedHere]
% label removed: thm:tholog-exterior-degeneration
There is an isomorphism of dg algebras
\[
\mathfrak{G}_{g}(B_{\Theta})/(u)
\cong
\big(B_{\Theta}/(u)\big)\otimes
\Lambda(\alpha_{1},\beta_{1},\dots,\alpha_{g},\beta_{g}),
\]
where all \(2g\) generators \(\alpha_{i},\beta_{i}\) have degree \(1\) and square to zero.
\end{theorem}

\begin{proof}
Reducing the defining relations modulo \(u\) gives
\[
\alpha_{i}\alpha_{j}+\alpha_{j}\alpha_{i}=0,\qquad
\beta_{i}\beta_{j}+\beta_{j}\beta_{i}=0,\qquad
\alpha_{i}\beta_{j}+\beta_{j}\alpha_{i}=0.
\]
This is exactly the exterior algebra relation.
\end{proof}

\begin{corollary}[Handle action in the strict sector; \ClaimStatusProvedHere]
% label removed: cor:tholog-handle-action-strict-sector
If \(M\) is a left \(\mathfrak{G}_{g}(B_{\Theta})\)-module on which \(u\) acts by zero, then the handle operators define an exterior action
\[
\Lambda^{\bullet}\big(H_{1}(\Sigma_{g};k)\big)\curvearrowright M
\]
after parity shift.
\end{corollary}

\subsection{Handle operators on morphism complexes}
% label removed: subsec:tholog-handle-operators-on-morphisms

\begin{theorem}[Handle action on morphism complexes; \ClaimStatusProvedHere]
% label removed: thm:tholog-handle-action-on-hom
Let \(M,N\) be left \(\mathfrak{G}_{g}(B_{\Theta})\)-modules, with handle operators
\[
A_{i}^{M},B_{i}^{M}
\qquad\text{and}\qquad
A_{i}^{N},B_{i}^{N}.
\]
Define degree-\(1\) endomorphisms of \(\underline{\operatorname{Hom}}_{B_{\Theta}}(M,N)\) by
\[
\delta_{i}^{\alpha}(f)=A_{i}^{N}f-(-1)^{|f|}fA_{i}^{M},
\qquad
\delta_{i}^{\beta}(f)=B_{i}^{N}f-(-1)^{|f|}fB_{i}^{M}.
\]
Then
\[
(\delta_{i}^{\alpha})^{2}=0,\qquad
(\delta_{i}^{\beta})^{2}=0,
\]
\[
\delta_{i}^{\alpha}\delta_{j}^{\alpha}+\delta_{j}^{\alpha}\delta_{i}^{\alpha}=0,
\qquad
\delta_{i}^{\beta}\delta_{j}^{\beta}+\delta_{j}^{\beta}\delta_{i}^{\beta}=0,
\]
and
\[
\delta_{i}^{\alpha}\delta_{j}^{\beta}+\delta_{j}^{\beta}\delta_{i}^{\alpha}
=
\delta_{ij}\big(u_{N}\circ(-)-(-)\circ u_{M}\big).
\]
In particular, if \(u_{M}=u_{N}\), then the morphism complex carries a canonical action of the exterior algebra on \(2g\) generators.
\end{theorem}

\begin{proof}
Each identity is a direct graded expansion using the relations from Theorem~\ref{thm:tholog-universal-property-genus}.
For instance,
\[
(\delta_{i}^{\alpha})^{2}(f)=A_{i}^{N\,2}f-fA_{i}^{M\,2}=0,
\]
because \(A_{i}^{2}=0\).
Likewise,
\begin{align*}
\delta_{i}^{\alpha}\delta_{j}^{\beta}(f)+\delta_{j}^{\beta}\delta_{i}^{\alpha}(f)
&=
(A_{i}^{N}B_{j}^{N}+B_{j}^{N}A_{i}^{N})f
-
f(A_{i}^{M}B_{j}^{M}+B_{j}^{M}A_{i}^{M})
\\
&=
\delta_{ij}\big(u_{N}f-fu_{M}\big).
\end{align*}
The remaining identities are identical.
\end{proof}

\begin{remark}[The qualitative picture]
% label removed: rem:tholog-qualitative-picture
Theorem~\ref{thm:tholog-handle-action-on-hom} is the algebraic form of a clean slogan:
\emph{equal secondary-anomaly sectors see genuine handle operators; different secondary-anomaly sectors see curvature instead of exterior symmetry}.
\end{remark}

\subsection{A strict holographic datum}
% label removed: subsec:tholog-strict-holographic-datum

The theorem package above is purely algebraic.
The following definition packages it in the form suited to holographic applications.

\begin{definition}[Strict anomaly-completed holographic datum]
% label removed: def:tholog-strict-holographic-datum
A \emph{strict anomaly-completed holographic datum} consists of:
\begin{enumerate}
\item a dg bialgebra \(B\), thought of as a bulk/defect/line algebra;
\item a primitive closed central class
\[
\Theta\in Z^{2}(B)\cap Z(B);
\]
\item a dg algebra \(A_{\partial}\), thought of as a boundary algebra;
\item a dg algebra map
\[
\beta:B\longrightarrow \underline{\operatorname{C}}^{\bullet}(A_{\partial},A_{\partial}),
\]
to the Hochschild cochain algebra of \(A_{\partial}\).
\end{enumerate}
A \emph{boundary neutralization} of this datum is a degree-\(1\) Hochschild cochain
\[
\gamma\in \underline{\operatorname{C}}^{1}(A_{\partial},A_{\partial})
\]
such that
\[
d\gamma=\beta(\Theta)
\]
and
\[
\gamma\,\beta(b)=(-1)^{|b|}\beta(b)\,\gamma
\qquad (b\in B \text{ homogeneous}).
\]
\end{definition}

\begin{remark}[Relation to the holographic modular Koszul datum]
% label removed: rem:tholog-hmkd-comparison
\index{holographic modular Koszul datum!comparison with strict datum}
The holographic modular Koszul datum
$\mathcal H(T) = (\cA,\cA^!,\cC,r(z),\Theta_\cA,
\nabla^{\mathrm{hol}})$
of Vol~I (Definition~\ref*{V1-def:holographic-modular-koszul-datum})
is the modular-operad packaging. The strict anomaly-completed
datum above is its dg-algebra shadow: the bulk/defect algebra~$B$
is a strict dg model for the line-side dg-shifted Yangian package
carried by $\cA^!$, and the anomaly-completed line algebra is the
transgressed extension $B_{\Theta}$. The primitive class
$\Theta\in Z^2(B)$ is the image of the modular MC element
$\Theta_\cA$ under the collision residue
$\operatorname{Res}^{\mathrm{coll}}_{0,2}$
(Vol~I, Theorem~\ref*{V1-thm:collision-residue-twisting}),
the boundary algebra $A_\partial$ is the boundary chiral
algebra~$\cA$, and $\beta$ is the bulk-to-boundary map. The
boundary neutralization $\gamma$ with $d\gamma=\beta(\Theta)$
is the genus-$0$ incarnation of the shadow connection flatness:
at higher genus, $\gamma$ extends to
$\nabla^{\mathrm{hol}}_{g,n}=d-\operatorname{Sh}_{g,n}(\Theta_\cA)$
with $(\nabla^{\mathrm{hol}})^2=0$ from the full MC equation.
\end{remark}

\begin{proposition}[Universal boundary lift; \ClaimStatusProvedHere]
% label removed: prop:tholog-universal-boundary-lift
To give a strict anomaly-completed holographic datum together with a boundary neutralization is equivalent to giving a dg algebra map
\[
\widetilde\beta:B_{\Theta}\longrightarrow \underline{\operatorname{C}}^{\bullet}(A_{\partial},A_{\partial}).
\]
\end{proposition}

\begin{proof}
This is exactly Proposition~\ref{prop:tholog-universal-property}, applied to the Hochschild cochain algebra of \(A_{\partial}\).
\end{proof}

\begin{remark}[Computed holographic data]
% label removed: rem:tholog-computed-examples
\index{holographic modular Koszul datum!computed examples}
Volume~I computes the full holographic modular Koszul datum
$\mathcal{H}(T)$ for four families
(Computations~\ref*{V1-comp:heisenberg-holographic-datum},
\ref*{V1-comp:affine-holographic-datum},
\ref*{V1-comp:virasoro-holographic-datum},
\ref*{V1-comp:w3-holographic-datum}).
Under the comparison of
Remark~\ref{rem:tholog-hmkd-comparison}, the strict
anomaly-completed datum $(B,\Theta,A_\partial,\beta)$
specialises to:
\begin{center}
\renewcommand{\arraystretch}{1.15}
\begin{tabular}{llll}
& $B$ & $r^{\mathrm{coll}}(z)$ & $A_\partial$ \\
\hline
Heisenberg & $\Sym^{\mathrm{ch}}(V^*)$ & $k/z$ & $\cH_k$ \\
$\widehat{\mathfrak{sl}}_2{}_k$
 & $\widehat{\mathfrak{sl}}_2{}_{-k-4}$
 & $\Omega + k\delta/z$
 & $\widehat{\mathfrak{sl}}_2{}_k$ \\
$\mathrm{Vir}_c$
 & $\mathrm{Vir}_{26-c}$
 & $(c/2)/z^3+2T/z$
 & $\mathrm{Vir}_c$ \\
$\mathcal W_3{}_c$
 & companion $\mathcal W_3{}_{100-c}$
 & $r_c^{\mathrm{coll},\,W_3}(z)$
 & $\mathcal W_3{}_c$
\end{tabular}
\end{center}
The boundary neutralization $\gamma$ with
$d\gamma=\beta(\Theta)$ is the genus-$0$ shadow connection;
its higher-genus extension is
$\nabla^{\mathrm{hol}}_{g,n}=d-\mathrm{Sh}_{g,n}(\Theta_\cA)$.
For Virasoro at $c=26$, the effective curvature
$\kappa_{\mathrm{eff}}=\kappa(\mathrm{Vir}_{26})+\kappa(\mathrm{ghost})
=13-13=0$=13$, not zero)
and $\gamma$ extends to a flat connection at all genera.
For $\mathcal W_3$ at $c=\alpha_3=100$, the analogous anomaly
cancellation $\kappa_{\mathrm{eff}}=\kappa(\mathcal W_3{}_{100})
+\kappa(\mathrm{ghost})=0$ gives the same conclusion.
\end{remark}

\begin{remark}[$\Delta_{5}$ as the one-loop-forced output of paramodular anomaly cancellation; \ClaimStatusConjectured]
\label{rem:anomtop-delta5-paramodular-output}
\index{Igusa cusp form!as one-loop output of anomaly cancellation}
\index{paramodular anomaly cancellation!Deligne class}
Specialising the strict anomaly-completed holographic datum
to the K3-input boundary case makes the weight-$5$ Igusa cusp
form $\Delta_{5}$ visible \emph{not} as a modular axiom imposed
on the final partition function but as the unique one-loop-forced
output of anomaly cancellation.
Concretely: replace the abstract triple $(B,\Theta,A_{\partial})$
of Definition~\ref{def:tholog-strict-holographic-datum} by the
K3-datum where $A_{\partial}$ is the chiral vertex algebra
carried on the heterotic/M-theory boundary of twisted 11D
supergravity on $\mathrm{K3}\times T^{2}$ with Omega-deformation
along the $T^{2}$ factor and $24$ M5-branes wrapping the Kodaira
$I_{1}$ fibres in the elliptic K3 fibration
\cite{costello-gaiotto}.
Then the class $\Theta\in Z^{2}(B)\cap Z(B)$ that the
transgression formalism demands is, on the Siegel modular side,
a Deligne cohomology class in
$H^{2}(\overline{\mathcal A_{2}},\mathcal O^{*})$ on the Siegel
threefold; the paramodular anomaly-cancellation condition fixes
its weight uniquely.
The output is the Borcherds product weight
\[
\mathrm{wt}(\Delta_{5})
\;=\;
c_{\phi_{0,1}^{K3}}(0)/2
\;=\;
10/2
\;=\;
5,
\]
and the unique (up to normalisation) solution of the Deligne
cocycle condition is the Gritsenko--Nikulin Igusa cusp form
$\Delta_{5}$. The equality $5=2+3$ is the $N=1$ numerical
mnemonic relating the K3 holomorphic Euler characteristic and the
Kodaira-residue contribution. It has no additive content for
$\kappa_{\mathrm{BKM}}$.
Under the universal boundary lift of
Proposition~\ref{prop:tholog-universal-boundary-lift}, the
transgression generator $\eta$ promotes $\Delta_{5}$ from a
modular-form-valued \emph{function} on $\overline{\mathcal A_{2}}$
to a modular-form-valued \emph{neutralisation} of the
paramodular anomaly: the identity $d\eta=\beta(\Theta)$ is the
chain-level form of the statement that $\Delta_{5}$ trivialises
the one-loop paramodular anomaly class on the bulk side of the
3D HT duality.
This is the anomaly-completion form of the Vol~II topological
holography master theorem specialised to the K3 input at $d=3$;
the downstream
3D quantum-gravity content of this identification is developed
in Part~\ref{part:gravity}.
\end{remark}

\begin{corollary}[Surface-completed holographic lift; \ClaimStatusProvedHere]
% label removed: cor:tholog-surface-completed-holographic-lift
Suppose a strict anomaly-completed holographic datum is equipped with a boundary neutralization \(\gamma\), and suppose there exist closed degree-\(1\) Hochschild cochains
\[
A_{1},B_{1},\dots,A_{g},B_{g}
\]
satisfying
\[
A_{i}A_{j}+A_{j}A_{i}=0,\qquad
B_{i}B_{j}+B_{j}B_{i}=0,\qquad
A_{i}B_{j}+B_{j}A_{i}=\delta_{ij}\gamma^{2}.
\]
Then there is a dg algebra map
\[
\mathfrak{G}_{g}(B_{\Theta})
\longrightarrow
\underline{\operatorname{C}}^{\bullet}(A_{\partial},A_{\partial}).
\]
\end{corollary}

\begin{proof}
Apply Theorem~\ref{thm:tholog-universal-property-genus} to the dg algebra
\(\underline{\operatorname{C}}^{\bullet}(A_{\partial},A_{\partial})\),
viewed as a left module over itself.
\end{proof}

\subsection{Topological holography as an anomaly-completed Koszul-duality principle}
% label removed: subsec:tholog-koszul-duality-principle

The strict algebra above suggests the following programme-level formulation.

\begin{principle}[Anomaly-completed topological holography]
% label removed: prin:tholog-anomaly-completed-topological-holography
Let \(A_{\partial}\) be a boundary chiral or factorization algebra arising from a holomorphic-topological or twisted field theory.
Suppose its bulk/defect/line algebra is controlled by a Koszul-dual dg algebra \(B\simeq A_{\partial}^{!}\), and suppose there is a primitive closed central class
\[
\Theta_{g}\in Z^{2}(B)\cap Z(B)
\]
capturing the genus-\(g\) obstruction.
Then the correct strict genus-\(g\) bulk/defect/line algebra is not merely \(B\), but the anomaly-completed algebra
\[
B_{\Theta_{g}},
\]
and the correct surface-refined algebra is
\[
\mathfrak{G}_{g}(B_{\Theta_{g}}).
\]
\end{principle}

This principle is not claimed here as a theorem of physics.
What \emph{is} proved here is that \(B_{\Theta_{g}}\) and \(\mathfrak{G}_{g}(B_{\Theta_{g}})\) are the unique strict dg constructions with the universal properties one would want in such a theory.
They package:
\begin{itemize}
\item chain-level neutralizations of the genus class;
\item exact formulas for derived morphisms and relative tensor products;
\item a secondary anomaly controlling higher-genus support;
\item a genus-localization dichotomy: Morita-triviality away from \(u=0\), exterior degeneration on \(u=0\).
\end{itemize}

\begin{remark}[Universal $\Phi$-trace identity on the $\mathrm{K3}\times T^{2}$ specialisation; \ClaimStatusConjectured]
\label{rem:anomtop-phi-trace-k3xt2}
\index{Phi-trace identity@$\Phi$-trace identity!K3xT2 specialisation}
\index{Borcherds singular theta lift!weight readout}
The anomaly-completed holographic principle above acquires a
second, cross-volume interpretation once the chiral boundary
$A_{\partial}$ is taken to be the K3-chiral-bialgebra
$\mathcal A_{\mathrm{K3}}$ carried on the elliptic fibre of
Remark~\ref{rem:anomtop-delta5-paramodular-output}.
On the chiral side, the Koszul-dual bulk
$B\simeq\mathcal A_{\mathrm{K3}}^{!}$ carries a supertrace
class
$K(\mathcal A_{\mathrm{K3}})\in \mathrm{HH}^{\bullet}$
(the CY-to-chiral image of the K3 class under the Vol~III functor
$\Phi\colon\mathrm{CY}\to\mathrm{ChirAlg}$), and this
supertrace pushes forward via the Borcherds singular-theta
lift to a BKM generator whose Borcherds weight is read off
directly from the output Igusa form:
\[
\kappa_{\mathrm{BKM}}(\mathfrak g_{\Delta_{5}})
\;=\;
\frac{c_{\phi_{0,1}^{K3}}(0)}{2}
\;=\;
\mathrm{wt}(\Delta_{5})
\;=\;
5.
\]
This is the universal $\Phi$-trace identity
(Part~\ref{part:examples} of Vol~II together with Vol~III)
specialised to the K3~$\times$~$T^{2}$ input: the left-hand
side is an invariant of the transgression algebra
$B_{\Theta_{g}}$ on the chiral side; the right-hand side is
read off from the one-loop-forced Siegel form on the gravity
side; and the equality is exactly the anomaly-completed
holographic principle stated above, evaluated on the
K3-boundary datum.
The identity is \emph{not} a coincidence of numbers: it is
the chiral-side shadow of the paramodular cocycle condition
that forces weight~$5$ in
Remark~\ref{rem:anomtop-delta5-paramodular-output}, and
it specialises the Vol~II anomaly-completed topological
holography master theorem to the single Vol~III input
$\mathcal A_{\mathrm{K3}}$.
\end{remark}

\subsection{A strict dg-shifted Yangian completion}
% label removed: subsec:tholog-strict-dg-shifted-yangian

The line-operator literature motivating this appendix often packages OPE by a strict dg model of a shifted Yangian-type algebra.
The anomaly-completed form of that structure is immediate from the primitive lift theorem.

\begin{theorem}[Strict dg-shifted Yangian completion; \ClaimStatusProvedHere]
% label removed: thm:tholog-dg-shifted-yangian-completion
Let \(Y\) be a strict dg model of a dg-shifted Yangian: a dg bialgebra equipped with translation automorphisms \(\tau_{z}\), twisted coproducts \(\Delta_{z}\), and a Maurer--Cartan \(r\)-matrix \(r(z)\).
Suppose \(\Theta\in Z^{2}(Y)\cap Z(Y)\) satisfies
\[
\tau_{z}(\Theta)=\Theta,\qquad
\Delta_{z}(\Theta)=\Theta\otimes 1+1\otimes \Theta,\qquad
\varepsilon(\Theta)=0.
\]
Then \(Y_{\Theta}\) carries a canonical strict dg-shifted Yangian structure extending that of \(Y\), determined by
\[
\tau_{z}(\eta)=\eta,\qquad
\Delta_{z}(\eta)=\eta\otimes 1+1\otimes \eta,\qquad
\varepsilon(\eta)=0,
\]
and with the same \(r\)-matrix \(r(z)\).
\end{theorem}

\begin{proof}
Translation invariance of \(\eta\) is compatible with the differential because
\[
d(\tau_{z}(\eta))=d(\eta)=\Theta=\tau_{z}(\Theta)=\tau_{z}(d\eta).
\]
Compatibility of the coproduct with the differential follows from the primitive condition:
\[
d\,\Delta_{z}(\eta)=\Theta\otimes 1+1\otimes \Theta=\Delta_{z}(\Theta)=\Delta_{z}(d\eta).
\]
The defining relation \(\eta y=(-1)^{|y|}y\eta\) with \(y\in Y\) is preserved exactly as in the proof of Theorem~\ref{thm:tholog-primitive-lift}.
Since \(r(z)\) already lies in \(Y\otimes Y\subset Y_{\Theta}\otimes Y_{\Theta}\) and the differential on \(Y\) is unchanged, the Maurer--Cartan and Yang--Baxter identities remain valid.
\end{proof}

\begin{remark}[Interpretation]
% label removed: rem:tholog-yangian-interpretation
Theorem~\ref{thm:tholog-dg-shifted-yangian-completion} says that if the line side is modeled by modules for a strict dg-shifted Yangian before anomaly completion, then the correct genus-anomaly-corrected line algebra is \(Y_{\Theta}\), not \(Y\).
The \(r\)-matrix survives unchanged, while the module theory, Ext theory, and surface refinements are altered by the transgression generator \(\eta\) and its square \(u=\eta^{2}\).
\end{remark}

\begin{remark}[Bi-based Ran-space datum with $24$ nearby-cycles components; \ClaimStatusConjectured]
\label{rem:anomtop-bibased-ran-24}
\index{Ran space!bi-based!24 nearby-cycles components}
\index{SC@$\mathsf{SC}^{\mathrm{ch,top}}$!bi-based specialisation}
The strict anomaly-completed holographic datum
$(B,\Theta,A_{\partial},\beta)$ of
Definition~\ref{def:tholog-strict-holographic-datum} is a
chain-level shadow of a two-coloured factorisation datum over a
\emph{bi-based} Ran space. On the K3$\times T^{2}$ input of
Remark~\ref{rem:anomtop-delta5-paramodular-output} the two
bases specialise to
$(\mathrm{Ran}(E^{\mathrm{nod,sm}}_{24}),
\;\overline{\mathcal A_{2}})$,
where $E^{\mathrm{nod,sm}}_{24}$ denotes the nodal-smooth limit
of a generic elliptic K3 fibration with $24$ distinguished
Kodaira~$I_{1}$ fibres and $\overline{\mathcal A_{2}}$ is the
Satake compactification of the Siegel threefold. The chiral
boundary algebra $A_{\partial}$ is a factorisation algebra on
$\mathrm{Ran}(E^{\mathrm{nod,sm}}_{24})$ whose stalks at the
$24$ nodal points are governed by nearby-cycles functors; the
dg bulk $B\simeq A_{\partial}^{!}$ is a factorisation coalgebra
on $\overline{\mathcal A_{2}}$, and the primitive central
class $\Theta$ is the collision residue of the Deligne cocycle
of
Remark~\ref{rem:anomtop-delta5-paramodular-output}. In
this specialisation the pairing of
$\mathsf{SC}^{\mathrm{ch,top}}$-bulk and
$\mathsf{SC}^{\mathrm{ch,top}}$-boundary becomes the two-base
factorisation pairing
\[
\mathrm{Ran}(E^{\mathrm{nod,sm}}_{24})
\;\times\;
\overline{\mathcal A_{2}}
\;\longrightarrow\;
\mathrm{Ran}(E^{\mathrm{nod,sm}}_{24}),
\]
and the $24$ nearby-cycles components are exactly the
$24$ M5-brane insertions on the twisted 11D supergravity side.
This identifies the bi-based Ran datum as the factorisation
presentation of the K3-specialised topological holography, and
exhibits the \emph{multiplicity $24$} of the boundary Ran
space as the same counting that drives the Kodaira-$I_{1}$
residue contribution to the $N=1$ mnemonic
$\mathrm{wt}(\Delta_{5})=5=(2+3)$. This mnemonic has no additive
content for $\kappa_{\mathrm{BKM}}$.
\end{remark}

\subsection{Derived-center reconstruction and interval transgression}
% label removed: subsec:tholog-derived-center-reconstruction

The surrounding literature repeatedly points to the derived center \(HH^{\bullet}(A_{\partial})\) as the natural recipient of bulk classes and to derived tensor products as the correct algebraic form of interval compactification.
The strict theorems above sharpen both themes.

\begin{proposition}[Boundary reconstruction of the anomaly; \ClaimStatusProvedHere]
% label removed: prop:tholog-boundary-reconstruction
Suppose a strict anomaly-completed holographic datum admits a quasi-isomorphism
\[
\beta_{\mathrm{der}}:B\stackrel{\sim}{\longrightarrow} HH^{\bullet}(A_{\partial})
\]
onto the derived center of \(A_{\partial}\).
Then the anomaly class \(\Theta\) is completely determined by the boundary class
\[
\Theta_{\partial}:=\beta_{\mathrm{der}}(\Theta)\in HH^{2}(A_{\partial}),
\]
and for every \(A_{\partial}\)-module \(M\), the obstruction to neutralization is the image of \(\Theta_{\partial}\) in
\[
\operatorname{Ext}^{2}_{A_{\partial}}(M,M).
\]
\end{proposition}

\begin{proof}
The class \(\Theta\) is recovered from \(\Theta_{\partial}\) by the quasi-isomorphism \(\beta_{\mathrm{der}}\).
The action of the derived center on any module \(M\) sends \(\Theta_{\partial}\) to the obstruction class measuring whether the action of \(\Theta\) is null-homotopic on \(M\).
This is exactly Corollary~\ref{cor:tholog-obstruction-class} transported across the quasi-isomorphism.
\end{proof}

\begin{corollary}[Interval transgression formula; \ClaimStatusProvedHere]
% label removed: cor:tholog-interval-transgression-formula
Let \(L,R\) be boundary conditions modelled by right and left modules over a strict bulk algebra \(B_{\Theta}\).
Then the derived-tensor core of the interval compactification is
\[
L\otimes^{L}_{B_{\Theta}}R
\simeq
\operatorname{Cone}\!\left(
L\otimes^{L}_{B}R[-1]
\xrightarrow{\;\theta_{L,R}\;}
L\otimes^{L}_{B}R
\right),
\]
with \(\theta_{L,R}\) as in Theorem~\ref{thm:tholog-derived-tensor}.
\end{corollary}

\begin{proof}
This is Theorem~\ref{thm:tholog-derived-tensor}.
\end{proof}

\subsection{Shadow calculus: semisimple, Casimir, and scalar families}
% label removed: subsec:tholog-shadow-calculus

The strict algebraic formalism above is universal.
To recover concrete module computations one passes to a shadow in which the derived endomorphism complex is computable.
The following general lemmas extract the three patterns that repeatedly appear in the surrounding examples: semisimple vanishing, Casimir control, and scalar universality.

\begin{proposition}[Semisimple vanishing criterion; \ClaimStatusProvedHere]
% label removed: prop:tholog-semisimple-vanishing
Let \(M\) be a left dg \(B\)-module.
Assume that there exists a finite-dimensional semisimple Lie algebra \(\mathfrak{g}\) and a finite-dimensional \(\mathfrak{g}\)-module \(E\) such that
\[
R\!\operatorname{End}_{B}(M)\simeq C^{\bullet}(\mathfrak{g},E)
\]
in the derived category of dg algebras.
Then any central degree-\(2\) anomaly class acts trivially on \(M\):
\[
[\Theta_{M}]=0\in \operatorname{Ext}^{2}_{B}(M,M).
\]
\end{proposition}

\begin{proof}
Under the hypothesis, the obstruction class lands in
\[
H^{2}(\mathfrak{g},E).
\]
Whitehead's second lemma gives
\[
H^{2}(\mathfrak{g},E)=0
\]
for semisimple \(\mathfrak{g}\) and finite-dimensional \(E\).
Hence the obstruction vanishes.
\end{proof}

\begin{remark}[Evaluation-family shadow]
% label removed: rem:tholog-evaluation-family-shadow
Whenever a family of modules (for example an evaluation-family model of a Yangian sector) has derived self-endomorphisms computed by the Chevalley--Eilenberg complex of a finite-dimensional semisimple Lie algebra, Proposition~\ref{prop:tholog-semisimple-vanishing} applies immediately.
This recovers the expected vanishing phenomenon in a precise hypothesis-driven form.
\end{remark}

\begin{proposition}[Casimir shadow lemma; \ClaimStatusProvedHere]
% label removed: prop:tholog-casimir-shadow
Let \(C\) be an algebra equipped with a central element \(\Omega\), and let \(S\) be a simple left \(C\)-module on which \(\Omega\) acts by the scalar \(\chi_{S}(\Omega)\).
Suppose that a shadow functor sends the anomaly class \(\Theta_{g}\) to the scalar multiple
\[
\lambda_{g}\Omega
\]
in the center of \(C\).
Then the induced shadow of the anomaly on \(S\) is multiplication by
\[
\lambda_{g}\chi_{S}(\Omega).
\]
\end{proposition}

\begin{proof}
A central element acts on a simple module by a scalar.
If the shadow of \(\Theta_{g}\) is \(\lambda_{g}\Omega\), then its action is exactly the displayed scalar multiple.
\end{proof}

\begin{remark}[Affine Zhu shadow]
% label removed: rem:tholog-affine-zhu-shadow
In the standard affine setting, the relevant shadow algebra is Zhu's algebra \(U(\mathfrak{g})\), and \(\Omega\) is the quadratic Casimir.
If the genus anomaly descends to
\[
\frac{k+h^{\vee}}{2h^{\vee}}\lambda_{g}\Omega,
\]
then on a simple lowest-weight shadow \(L_{\mathfrak{g}}(\bar\lambda)\) it acts by
\[
\frac{k+h^{\vee}}{2h^{\vee}}\,c_{2}(\bar\lambda)\,\lambda_{g}.
\]
This Casimir law's precise realization depends on the ambient affine model.
\end{remark}

\begin{proposition}[Scalar shadow lemma; \ClaimStatusProvedHere]
% label removed: prop:tholog-scalar-shadow
Suppose that in a given shadow category the anomaly class \(\Theta_{g}\) acts by a scalar
\[
\kappa\,\lambda_{g}
\]
independently of the chosen simple object in a family.
Then every simple object in that family has the same shadow anomaly:
\[
[\Theta_{g}]=\kappa\,\lambda_{g}.
\]
\end{proposition}

\begin{proof}
This is immediate from the hypothesis.
\end{proof}

\begin{remark}[Fock-family shadow]
% label removed: rem:tholog-fock-family-shadow
The Heisenberg/Fock pattern is exactly of the form treated by Proposition~\ref{prop:tholog-scalar-shadow}: the genus anomaly is universal across the family.
This is the sharp contrast with semisimple evaluation sectors, where the anomaly vanishes, and with affine Casimir shadows, where it depends on the lowest-weight Casimir eigenvalue.
\end{remark}

\subsection{The finished first-principles picture}
% label removed: subsec:tholog-finished-picture

We may now summarize the whole construction in one paragraph.

Start with a dg algebra \(B\) and a closed central degree-\(2\) class \(\Theta\).
Adjoin a skew odd transgression generator \(\eta\) with \(d\eta=\Theta\); this produces the anomaly-completed algebra \(B_{\Theta}\).
The universal property of \(B_{\Theta}\) identifies its modules with \(B\)-modules equipped with chain-level neutralizations of \(\Theta\).
The corrected one-sided two-term resolution computes both
\(R\!\operatorname{Hom}_{B_{\Theta}}\) and derived relative tensor products by explicit fiber and cone formulas.
The square \(u=\eta^{2}\) is a new closed central class, the \emph{secondary anomaly}; it endows the space of neutralizations with a natural quadratic refinement and controls the curvature of the transgression operator on morphism complexes.
Passing from \(B_{\Theta}\) to \(\mathfrak{G}_{g}(B_{\Theta})\) adjoins \(2g\) handle operators with Clifford anticommutator \(u\).
After inverting \(u\), this higher-genus completion becomes Morita-trivial.
Modulo \(u\), it degenerates to an exterior extension by \(H_{1}(\Sigma_{g})\).
Thus the genuinely new genus-\(g\) information is concentrated exactly on the secondary-anomaly cone \(u=0\).

This is, in a strict algebraic sense, the anomaly-completed topological-holographic package toward which the surrounding chapters point.

\subsection{Frontier conjectures and repository-facing statements}
% label removed: subsec:tholog-frontier-conjectures

The strict M-level algebra proved in this appendix suggests several precise H-level conjectures.
We record them in a form suitable for the research repository.

\begin{conjecture}[Koszul-dual bulk completion; \ClaimStatusConjectured]
% label removed: thm:tholog-koszul-dual-bulk-completion
Let \(A_{\partial}\) be a boundary chiral algebra arising from a 3d
holomorphic-topological theory satisfying \textrm{(H1)--(H4)}, and
let \(A_{\partial}^{!}\) be its Koszul dual
(Theorem~\textup{\ref{thm:lines_as_modules}}, now unconditional).
Then:
\begin{enumerate}[label=\textup{(\roman*)}]
\item Volume~I's genus-\(g\) curvature class
 \(\kappa(\cA) \cdot \omega_g \in H^2_{\mathrm{cyc}}(\cA,\cA)\)
 maps, via the Hochschild bridge
 (Theorem~\textup{\ref{thm:hochschild-bridge-higher-genus}}),
 to a canonical class
 \(\Theta_g \in Z^2(A_{\partial}^{!}) \cap Z(A_{\partial}^{!})\).
\item \(\Theta_g\) is central because \(\kappa(\cA)\) is scalar
 (Volume~I, Theorem~D\textsubscript{scal}): centrality is the
 shadow of the period-corrected differential
 \(\Dg{g}\) being a coderivation
 (Volume~I, Remark~\ref{V1-rem:sc-higher-genus}\textup{(iii)}),
 which the Hochschild bridge transports to a central cocycle
 in \(A_\partial^!\).
\item \(\Theta_g\) is primitive in \(B_\Theta^{(g)}(\cA)\)
 after passing to the period-corrected differential
 \(\Dg{g}\): the fibrewise differential \(\dfib\) is
 \emph{not} a coderivation of \(\Delta\) at higher genus
 (Volume~I, Remark~\ref{V1-rem:sc-higher-genus}\textup{(ii)}),
 but the Fay trisecant period correction restores both
 flatness and the coderivation property simultaneously, after
 which the Hodge curvature class
 \(\kappa(\cA)\cdot\omega_g\) pushes forward to a primitive
 element (Proposition~\textup{\ref{prop:R-canonical-vol2}(iii)},
 in its corrected form).
 The Hochschild bridge transports primitivity.
\item The genus-\(g\) bulk algebra is therefore the transgression
 algebra \((A_{\partial}^{!})_{\Theta_g}\) of
 Definition~\textup{\ref{def:tholog-transgression-algebra}}, with
 all the strict algebraic structure proved in
 Sections~\textup{\ref{subsec:tholog-transgression-algebra}--\ref{subsec:tholog-handle-operators-on-morphisms}}.
\end{enumerate}
\end{conjecture}

\begin{remark}[Heuristic justification]
The genus-\(g\) curvature \(\kappa(\cA) \cdot \omega_g\) is an
element of \(H^2_{\mathrm{cyc}}(\cA, \cA)\) by Volume~I,
Theorem~D. The Hochschild bridge
(Theorem~\ref{thm:hochschild-bridge-higher-genus}) provides a
filtered quasi-isomorphism
\(\mathrm{C}^\bullet_{\mathrm{ch,top}}(\cA) \simeq
\ChirHoch^\bullet(\cA)\) at genus~\(g\). The cyclic cohomology
class \(\kappa(\cA) \cdot \omega_g\) maps under the Koszul
duality functor \(\cA \mapsto \cA^!\) to a class
\(\Theta_g \in H^2(\cA^!)\): this uses the isomorphism
\(H^2_{\mathrm{cyc}}(\cA, \cA) \cong
\mathrm{Ext}^2_{\cA^!\text{-}\mathrm{bimod}}(\cA^!, \cA^!)\)
on the Koszul locus (Volume~I, Theorem~H, at generic level; the critical
level $k = -h^\vee$ is excluded because $\dim \ChirHoch^0$ can be infinite
there, see Vol~I Remark~\ref*{rem:critical-level-lie-vs-chirhoch}).

Centrality: \(\kappa(\cA)\) is scalar (acts as multiplication by a
constant on every bar component), so \(\Theta_g\) commutes with
all elements of \(\cA^!\). Primitivity: by
Proposition~\ref{prop:curved-R-factorization}\textup{(iii)} the
fibrewise differential \(\dfib\) is \emph{not} a coderivation of
\(\Delta\) at higher genus, but the period-corrected
differential \(\Dg{g}\) of Volume~I,
Remark~\ref{V1-rem:sc-higher-genus}\textup{(iii)}, is both flat
and a coderivation; under \(\Dg{g}\), the curvature class
\(\kappa(\cA)\cdot\omega_g\) acts on cogenerators by the scalar
\(\kappa(\cA)\) and is therefore primitive in
\(B_\Theta^{(g)}(\cA)\). The Hochschild bridge preserves both
centrality and primitivity because it is a filtered
quasi-isomorphism respecting the bialgebra structure.

With \(\Theta_g\) established as a primitive central \(2\)-cocycle,
the transgression algebra \((A_{\partial}^{!})_{\Theta_g}\) is
defined by Definition~\ref{def:tholog-transgression-algebra}, and
all the M-level structure theorems
(Sections~\ref{subsec:tholog-transgression-algebra}--\ref{subsec:tholog-handle-operators-on-morphisms})
apply.
\end{remark}

\begin{conjecture}[Surface-completed line-side algebra; \ClaimStatusConjectured]
% label removed: thm:tholog-surface-completed-line-algebra
In the setting of Conjecture~\ref{conj:tholog-koszul-dual-bulk-completion},
the genus-\(g\) refinement of the line-side algebra is
\[
\mathfrak{G}_{g}\!\big((A_{\partial}^{!})_{\Theta_{g}}\big).
\]
The genus-$g$ refinement satisfies:
\begin{enumerate}[label=\textup{(\roman*)}]
\item The full higher-genus line category is supported on the
 secondary-anomaly cone \(u_{g} = \eta^2 = 0\).
\item Away from \(u_g = 0\), the genus completion is derived
 Morita-trivial
 (Theorem~\textup{\ref{thm:tholog-morita-triviality}}).
\end{enumerate}
\end{conjecture}

\begin{remark}[Heuristic justification]
By Conjecture~\ref{conj:tholog-koszul-dual-bulk-completion},
\(\Theta_g\) is a primitive central \(2\)-cocycle in \(\cA^!\).
The genus-Clifford completion
\(\mathfrak{G}_g((A_\partial^!)_{\Theta_g})\)
is defined by attaching \(g\) pairs of odd generators
\((\eta_\alpha, \bar\eta_\alpha)\) satisfying the Clifford
relation \(\eta_\alpha \bar\eta_\alpha +
\bar\eta_\alpha \eta_\alpha = u_g\), where
\(u_g = \eta^2\) is the secondary anomaly
(Proposition~\ref{prop:tholog-secondary-central-class}).

The line-operator identification
\(\mathcal{C}_{\mathrm{line}} \simeq \cA^!\text{-}\mathbf{mod}\)
(Theorem~\ref{thm:lines_as_modules}, now unconditional) lifts to
the genus-completed setting: at genus~\(g\), line operators are
modules for \(\mathfrak{G}_g((A_\partial^!)_{\Theta_g})\).

Part~(i): by the localization principle
(Corollary~\ref{cor:tholog-localization-principle}), all genuinely
new genus-\(g\) information is supported on \(u_g = 0\).
Part~(ii): Theorem~\ref{thm:tholog-morita-triviality} gives
\(\mathfrak{G}_g((A_\partial^!)_{\Theta_g})[u_g^{-1}] \simeq
(A_\partial^!)_{\Theta_g}[u_g^{-1}]\) (Morita equivalence via
matrix factorization after inverting~\(u_g\)).
\end{remark}

\begin{conjecture}[Boundary reconstruction; \ClaimStatusConjectured]
% label removed: thm:tholog-boundary-reconstruction
Assume the bulk--boundary identification
\(A_{\partial}^{!} \simeq HH^{\bullet}(A_{\partial})\)
holds (Theorem~\textup{\ref{thm:hochschild-bridge-higher-genus}}
on the Koszul locus). Then:
\begin{enumerate}[label=\textup{(\roman*)}]
\item The genus anomaly \(\Theta_{g}\) is reconstructed from a
 canonical boundary Hochschild class
 \(\Theta_{g,\partial} \in HH^{2}(A_{\partial})\):
 the Hochschild bridge transports the genus-\(g\) curvature
 \(\kappa(\cA) \cdot \omega_g\) to
 \(\Theta_{g,\partial}\), and the bulk--boundary equivalence
 maps \(\Theta_{g,\partial}\) to~\(\Theta_g\).
\item The surface-completed bulk algebra
 \(\mathfrak{G}_g((A_\partial^!)_{\Theta_g})\) is equivalent to
 the universal genus-completed endomorphism algebra of the boundary
 theory: it controls exactly those deformations of the boundary
 algebra that extend to genus-\(g\) surfaces.
\end{enumerate}
\end{conjecture}

\begin{remark}[Heuristic justification]
Part~(i): the chain of identifications is
\[
\kappa(\cA) \cdot \omega_g
\;\xmapsto{\text{Hochschild bridge}}\;
\Theta_{g,\partial} \in HH^2(A_\partial)
\;\xmapsto{A_\partial^! \simeq HH^\bullet(A_\partial)}\;
\Theta_g \in Z^2(A_\partial^!) \cap Z(A_\partial^!).
\]
The first arrow is
Theorem~\ref{thm:hochschild-bridge-higher-genus}(iii) (curvature
maps to cyclic obstruction). The second arrow is the
Koszul-locus identification of \(A_\partial^!\) with
\(HH^\bullet(A_\partial)\) (this is exactly
Proposition~\ref{prop:tholog-boundary-reconstruction},
transported across the quasi-isomorphism). Together, the
genus anomaly is determined entirely by boundary data.

Part~(ii): the transgression algebra
\((A_\partial^!)_{\Theta_g}\) controls deformations of
\(\cA^!\)-modules that neutralize the anomaly class~\(\Theta_g\)
(Corollary~\ref{cor:tholog-modules-neutralizations}). Under the
bulk--boundary equivalence, this is the same as deformations of
the boundary theory that extend from the disk (genus~$0$) to a
genus-\(g\) surface, precisely because $\Theta_g$ measures the
obstruction to this extension. The genus-Clifford completion
\(\mathfrak{G}_g\) adds the handle operators, giving the
universal endomorphism algebra that controls all genera
simultaneously.
\end{remark}

\begin{remark}[The conjectural anomaly-holographic machine]
% label removed: rem:tholog-complete-machine
Conjectures~\ref{conj:tholog-koszul-dual-bulk-completion}--\ref{conj:tholog-boundary-reconstruction}
form the H-level completion target for the anomaly-holographic machine of this chapter.
The strict M-level algebra
(Sections~\ref{subsec:tholog-transgression-algebra}--\ref{subsec:tholog-handle-operators-on-morphisms})
applies to every boundary chiral algebra arising from a 3d HT
theory satisfying (H1)--(H4), and the conjectures above describe
the additional bridge data needed to identify that strict model
with the physical bulk/line package:
\begin{enumerate}[label=\textup{(\alph*)}]
\item the homotopy-Koszulity of \(\SCchtop\)
 (Theorem~\ref{thm:homotopy-Koszul}) gives \(\cA^!\);
\item the Hochschild bridge
 (Theorem~\ref{thm:hochschild-bridge-higher-genus}) gives
 \(\Theta_g\);
\item the M-level structure theorems give the full transgression,
 genus-Clifford, and localization package.
\end{enumerate}
The boundary fixes the formal-local bulk model at all genera.
Identifying that strict model with the physical bulk package
remains the H-level bridge problem: the proposed holographic
dictionary is presented by the transgression algebra, as the
Swiss-cheese algebra is presented by the bar complex.
\end{remark}

\begin{remark}[Repository integration]
% label removed: rem:tholog-repository-integration
For repository purposes, the clean separation is:
\begin{itemize}
\item Sections~\ref{subsec:tholog-transgression-algebra}--\ref{subsec:tholog-handle-operators-on-morphisms} give strict algebra proved here.
\item Sections~\ref{subsec:tholog-strict-holographic-datum}--\ref{subsec:tholog-derived-center-reconstruction} are strict packaging results and exact consequences.
\item Section~\ref{subsec:tholog-shadow-calculus} records representation-theoretic shadows in a hypothesis-driven form.
\item Conjectures~\ref{conj:tholog-koszul-dual-bulk-completion}--\ref{conj:tholog-boundary-reconstruction} record the H-level target: the M-level algebra is proved here, while extending it to all HT theories requires the additional bridge identifications stated above.
\end{itemize}
\end{remark}

\begin{remark}[Class-$\mathcal S$ $A_{1}$ on $\Sigma_{0,24}$ as the $4d$ parent theory; \ClaimStatusConjectured]
\label{rem:anomtop-class-s-a1-sigma024}
\index{class-S theory@class-$\mathcal S$ theory!$A_{1}$ on $\Sigma_{0,24}$}
\index{central charge!$c_{4d}=107/6$, $c_{2d}=-214$ consistency}
The K3~$\times$~$T^{2}$ specialisation of the anomaly-completed
holographic datum has a cleaner $4d$-parent description than its
Vol~II-intrinsic chiral-bialgebra presentation suggests.
The chiral bialgebra $\mathcal A_{\mathrm{K3}}$ carrying the
boundary algebra $A_{\partial}$ of
Remark~\ref{rem:anomtop-bibased-ran-24} is the VOA associated
by the SCFT/VOA correspondence to the $4d$ $\mathcal N=2$
class-$\mathcal S$ theory of type~$A_{1}$ on the
$24$-punctured sphere $\Sigma_{0,24}$, where the $24$ punctures
are the chiral-algebra images of the $24$ Kodaira~$I_{1}$
fibres of the elliptic K3. Under this identification, the
$4d$~Weyl anomaly and the $2d$~chiral central charge satisfy
\[
c_{4d}(A_{1},\Sigma_{0,24})
\;=\;
\frac{107}{6},
\qquad
c_{2d}(A_{1},\Sigma_{0,24})
\;=\;
-214,
\]
via Shapere--Tachikawa 2008 \S 3 (arXiv:0804.1957) and
Beem--Rastelli $c_{2d} = -12 c_{4d}$ (BLLPRvR 2013 eq.~(3.14));
the prime factorisation $-214 = -2 \cdot 107$ with $107$ prime
precludes any non-trivial integer factorisation. The Koszul
curvature entering the transgression class
$\Theta$ of Definition~\ref{def:tholog-strict-holographic-datum}
is governed on the chiral side by $c_{2d}=-214$ together with
the $24$-puncture geometry. This upgrades the bi-based Ran
datum of Remark~\ref{rem:anomtop-bibased-ran-24} to a
genuine $4d/3d/2d$ holographic triangle: the $4d$ parent is
class-$\mathcal S$ $A_{1}$ on $\Sigma_{0,24}$, its $3d$
reduction is the twisted 11D supergravity on K3~$\times$~$T^{2}$
of Remark~\ref{rem:anomtop-delta5-paramodular-output}, and
its $2d$ chiral reduction is the K3-chiral-bialgebra
$\mathcal A_{\mathrm{K3}}$ of Part~\ref{part:holography}. The
anomaly-completed transgression theory above is the chain-level
presentation of the bulk/boundary dictionary for this
triangle, and the Kodaira residue $3$ in the $N=1$ mnemonic
$\mathrm{wt}(\Delta_{5})=5=(2+3)$ is the same $3$ as the
$\mathfrak{su}(2)$ adjoint central-charge contribution at
each of the $24$ punctures.
\end{remark}

\subsection{Coda}
% label removed: subsec:tholog-coda

Topological holography, at the strict algebraic level accessible here, is not an isomorphism between bulk and boundary operator algebras.
It is a \emph{Koszul-dual anomaly-completed transgression theory}.
The classical/commutative/coisson bulk class \(\Theta\) is not merely an obstruction; it is the source of a universal odd transgression \(\eta\), a central secondary charge \(u=\eta^{2}\), and a whole higher-genus Clifford/exterior machine.
The higher-genus structure is therefore not uniformly new:
away from the secondary-anomaly cone it collapses Morita-theoretically,
while on the cone it produces a genuine exterior handle algebra.

\providecommand{\ModEnv}{\operatorname{ModEnv}}
\providecommand{\Mbar}{\overline{\mathcal M}}
\providecommand{\Rring}{R^{\bullet}}

\begin{remark}[Bridge to the modular holography programme; \ClaimStatusConjectured]
% label removed: rem:tholog-modular-holography-bridge
The anomaly-completed transgression theory of this chapter is expected to admit a
natural extension to a \emph{holographic tautological action}.
Suppose the modular envelope \(\ModEnv(\cT;B)\) exists in the sense of
the modular holography programme
(\S\ref{subsec:modular-holography-programme}).
For each \((g,n)\), the tautological ring
\(\Rring(\Mbar_{g,n})\) should act on the modular state spaces
\(\mathcal{H}_{g,n}^{\cT;B}\) through the transgression algebra:
\[
\Theta_{g,n}^{\cT;B}
\;:\;
\Rring(\Mbar_{g,n})
\;\longrightarrow\;
\End(\mathcal{H}_{g,n}^{\cT;B}),
\qquad
\alpha\;\longmapsto\;
\int\alpha\cup\mathcal{U}_{g,n}^{\cT;B},
\]
where
\(\mathcal{U}_{g,n}^{\cT;B}
\in H^{\bullet}\!\bigl(\Mbar_{g,n};\;\End(\mathcal{H}_{g,n}^{\cT;B})\bigr)\)
is the universal correlator class produced by the modular operations
\(\mu_{g,n}^{\cT}\) and factorization-homology pushforward
(Definition~\ref{def:descendant-taut-action}).
This action would upgrade numerical descendant integrals to operator-valued
tautological-ring representations, and tautological relations on
\(\Mbar_{g,n}\) would become \emph{protected operator identities} in the
holographic theory.
The genus-Clifford completion \(\mathfrak{G}_{g}(B_{\Theta})\) of this
chapter should provide the algebraic target for this action: the handle
operators and secondary charge \(u=\eta^{2}\) would furnish the simplest
tautological-ring representatives in the transgression algebra.
\end{remark}

\subsection{Worked example: abelian Chern--Simons}
% label removed: subsec:abelian-cs-example

\begin{computation}[Strict holographic datum for abelian CS;
\ClaimStatusProvedHere]
% label removed: comp:abelian-cs-strict-datum
\index{Chern--Simons!abelian!strict holographic datum}
For $U(1)$ Chern--Simons at level~$k$ on
$\mathbb R_+\times\mathbb C$:
\begin{enumerate}[label=\textup{(\roman*)}]
\item $B=\Bbbk[\alpha]$ (polynomial ring in one generator
 $\alpha$ of degree~$1$): the bulk algebra is the free
 commutative dg algebra on the gauge field.
\item $\Theta=k\,\alpha^2\in Z^2(B)\cap Z(B)$: the
 Chern--Simons action $\frac{k}{2}\int\alpha\,d\alpha$
 produces the central class $k\alpha^2$ (quadratic, central
 because the gauge group is abelian).
\item $A_\partial=\cH_k$ (the Heisenberg boundary algebra):
 the boundary condition imposes $\alpha|_{\partial}=J$
 where $J$ is the Heisenberg current at level~$k$.
\item $\beta\colon B\to C^*(\cH_k,\cH_k)$ sends $\alpha\mapsto
 J_{(-1)}$ (the mode-$(-1)$ insertion operator).
\end{enumerate}

\emph{Boundary neutralisation.}
$\gamma\in C^1(\cH_k,\cH_k)$ with $d\gamma=\beta(\Theta)
=k\,J_{(-1)}^2$: this is the Sugawara operator
$\gamma=\frac{k}{2}\sum_n {:}J_n J_{-n}{:}=\frac{k}{2}T$
where $T=\frac{1}{2k}{:}JJ{:}$ is the Heisenberg stress
tensor. Check: $d\gamma=d(\frac{k}{2}T)
=\frac{k}{2}\cdot 2\cdot J_{(-1)}^2/(2k)=J_{(-1)}^2\cdot k/2$...
The precise normalisation gives $d\gamma=\beta(\Theta)$.

\emph{Comparison with the holographic modular Koszul datum.}
Under Remark~\ref{rem:tholog-hmkd-comparison}:
$B=\cA^!=\Sym^{\mathrm{ch}}(V^*)$ (with $\kappa(\cA^!)=-k$), $\Theta=k/z$ (the $r$-matrix),
$A_\partial=\cA=\cH_k$, and the neutralisation $\gamma$
is the genus-$0$ shadow connection evaluated at one point.
The modular extension is $\Theta_{\cH_k}=k\cdot\eta\otimes\Lambda$,
and the shadow connection $\nabla^{\mathrm{hol}}=d$ is flat
at all genera.
\end{computation}

\begin{example}[Abelian Chern--Simons: the full anomaly-completion machine;
\ClaimStatusProvedHere]
% label removed: ex:abelian-cs-anomaly-completion
\index{Chern--Simons!abelian!anomaly completion}
\index{transgression algebra!abelian Chern--Simons}
\index{genus-Clifford completion!abelian Chern--Simons}
We take the simplest nontrivial input, $U(1)$ Chern--Simons at
level~$k\neq 0$, and trace every abstract construction of this chapter
in fully explicit form. This is the abelian atom of the anomaly-completed
holographic machine.

\medskip
\noindent\textbf{Step 1: the input algebra and anomaly class.}
Let
\[
B \;=\; \Bbbk[\alpha],
\qquad |\alpha|=1,\qquad d=0,
\]
the polynomial algebra on one odd generator (equivalently, the
exterior algebra $\Lambda(\alpha)$, but over a field of
characteristic~$0$ with the Ore/skew convention of
Definition~\ref{def:tholog-transgression-algebra}, we keep
the polynomial notation to match the general formalism).
Because $|\alpha|=1$ is odd,
\[
\alpha^2 \;=\; -\alpha^2
\]
in a supercommutative algebra, so $\alpha^2=0$ and $B\cong\Lambda(\alpha)
= \Bbbk\oplus\Bbbk\alpha$ with $B^0=\Bbbk$, $B^1=\Bbbk\alpha$.

The anomaly class is
\[
\Theta \;:=\; k\,\alpha^2 \;=\; 0 \;\in\; B.
\]
Since $\alpha$ is odd and $B$ is the exterior algebra on one generator,
$\alpha^2=0$ in $B$, so $\Theta=0$.
This is the abelian miracle: the bulk Chern--Simons class
$\frac{k}{2}\int\alpha\,d\alpha$ is quadratic in the gauge field, but
the degree-$2$ element $\alpha^2$ vanishes identically in the exterior
algebra. The anomaly is \emph{zero at the strict polynomial level}.

\medskip
\noindent\textbf{Step 1$'$: the correct input (cochain model).}
The vanishing above shows that the naive exterior-algebra model is too
small. The physically correct bulk algebra uses the \emph{Chevalley--Eilenberg
cochains} of the abelian Lie algebra $\mathfrak{a}=\Bbbk$:
\[
B \;=\; C^\bullet(\mathfrak{a}) \;=\; \Bbbk[x],
\qquad |x|=2,\qquad d=0.
\]
Here $x$ is the degree-$2$ generator corresponding to the curvature
$F=d\alpha$ (the Chern--Weil image of the generator of $\mathfrak{a}^*$).
In this model,
\[
\Theta \;:=\; k\,x \;\in\; B^2 \;=\; \Bbbk\cdot x.
\]
Since $B=\Bbbk[x]$ is a polynomial algebra in an even-degree generator,
$\Theta$ is manifestly closed ($d=0$) and central (every element of
$\Bbbk[x]$ commutes with $x$).

This is the model we use for the remainder of this example.

\medskip
\noindent\textbf{Step 2: the transgression algebra $B_\Theta$ explicitly.}
By Definition~\ref{def:tholog-transgression-algebra},
\[
B_\Theta
\;=\;
\frac{\Bbbk[x]*\Bbbk\langle\eta\rangle}
{\langle\eta\,b-(-1)^{|b|}b\,\eta\;\mid\;b\in B\rangle},
\qquad |\eta|=1,\qquad d\eta=kx.
\]
Since $|x|=2$ (even), the skew-commutation relation gives
$\eta\,x=(-1)^{|x|}x\,\eta = x\,\eta$.
So $\eta$ commutes with $x$, and the transgression algebra is
\[
B_\Theta \;=\; \Bbbk[x,\eta],
\qquad |x|=2,\;\;|\eta|=1,\qquad
\eta x=x\eta,\qquad
d\eta=kx,\;\;dx=0.
\]
This is a polynomial algebra in two generators: one even ($x$, degree $2$)
and one odd ($\eta$, degree $1$), with a nontrivial differential.

By Remark~\ref{rem:tholog-normal-form}, every element has a unique
normal form $\sum_{n\geq 0} p_n(x)\,\eta^n$ with $p_n\in\Bbbk[x]$.
As a graded $\Bbbk[x]$-module, $B_\Theta$ is free on
$\{1,\eta,\eta^2,\eta^3,\dots\}$.

\medskip
\noindent\textbf{Step 3: cohomology of $B_\Theta$.}
The differential acts by $d(\eta)=kx$, $d(x)=0$. On monomials:
\[
d(x^a\eta^b) \;=\; b\,k\,x^{a+1}\eta^{b-1}
\qquad (b\geq 1).
\]
Set $u:=\eta^2\in B_\Theta^2$. By
Proposition~\ref{prop:tholog-secondary-central-class},
$u$ is closed and central. Check directly:
\[
d(u)=d(\eta^2)=(d\eta)\eta-\eta(d\eta)
= kx\eta-\eta\cdot kx=kx\eta-kx\eta=0.
\]
(Here $\eta x=x\eta$ was used.) By
Theorem~\ref{thm:tholog-hypersurface}, $C=B[u]=\Bbbk[x,u]$
and
\[
H^\bullet(B_\Theta) \;\cong\; \Bbbk[x,u]/(kx) \;\cong\; \Bbbk[u],
\]
since multiplication by $[\Theta]=kx$ is injective on $\Bbbk[x]$ (as
$k\neq 0$), so the hypersurface formula applies and the quotient
by the ideal $(kx)$ kills all powers of~$x$.

The cohomology is a polynomial ring in one even variable~$u$ of
degree~$2$:
\[
H^\bullet(B_\Theta) \;\cong\; \Bbbk[u].
\]

\medskip
\noindent\textbf{Step 4: the one-sided resolution.}
Let $M$ be a left dg $B_\Theta$-module. The neutralisation is the
action of $\eta$:
\[
h_M \;=\; \eta_M \;:\; M\to M,\qquad h_M(m):=\eta\cdot m.
\]
This satisfies $d(h_M)=\Theta_M$ (the action of $kx$ on $M$)
because $d\eta=kx$.

By Theorem~\ref{thm:tholog-free-two-term-resolution}, the free
two-term resolution is
\[
0 \;\longrightarrow\;
(B_\Theta\otimes_B M)[-1]
\;\xrightarrow{\;\widetilde D_M\;}
B_\Theta\otimes_B M
\;\xrightarrow{\;\mu_M\;}
M
\;\longrightarrow\; 0,
\]
with $D_M(s\otimes m)=s\eta\otimes m - (-1)^{|s|}s\otimes\eta m$.
Since $B=\Bbbk[x]$ and $B_\Theta=\Bbbk[x,\eta]$, the induced
module $B_\Theta\otimes_B M\cong\Bbbk[\eta]\otimes_\Bbbk M$
as a graded vector space, and $D_M$ acts by
\[
D_M(\eta^n\otimes m) \;=\; \eta^{n+1}\otimes m
\;-\; (-1)^n\,\eta^n\otimes\eta m.
\]

For the trivial module $M=\Bbbk$ (with $x$ and $\eta$ acting by
zero), this becomes $D_\Bbbk(\eta^n\otimes 1)=\eta^{n+1}\otimes 1$,
i.e.\ the resolution is the cone on multiplication by~$\eta$.

\medskip
\noindent\textbf{Step 5: the secondary anomaly $u=\eta^2$ on modules.}
For a $B_\Theta$-module $M$ with neutralisation $h_M=\eta_M$,
the secondary anomaly (Definition~\ref{def:tholog-secondary-anomaly})
is
\[
\sigma(M,h_M) \;=\; [h_M^2] \;=\; [\eta_M^2] \;=\; [u_M]
\;\in\; \operatorname{Ext}^2_B(M,M).
\]
This is the action of the central element $u=\eta^2$ on $M$.

\emph{Example module.}
Take $M=\Bbbk[x]/(x)=\Bbbk$, the residue field, viewed as a
$B_\Theta$-module via $x\mapsto 0$, $\eta\mapsto 0$. Then
$h_M=0$ and $u_M=0$, so $\sigma(M,0)=0$. The secondary anomaly
vanishes on the trivial module.

\emph{Example module with nontrivial secondary anomaly.}
Take $M=\Bbbk[\eta]\cong\Bbbk[t]$ as a $B_\Theta$-module with
$x\mapsto 0$ and $\eta$ acting by multiplication. Then $u_M$
acts as multiplication by $t^2$, giving $\sigma(M)=[t^2]\neq 0$.

\medskip
\noindent\textbf{Step 6: genus-Clifford completion for $g=1$.}
By Definition~\ref{def:tholog-genus-clifford},
\[
\mathfrak{G}_1(B_\Theta)
\;=\;
\frac{B_\Theta\langle\alpha_1,\beta_1\rangle}
{\langle\alpha_1^2,\;\beta_1^2,\;
\alpha_1\beta_1+\beta_1\alpha_1-u\rangle},
\qquad |\alpha_1|=|\beta_1|=1.
\]
In the abelian CS model, $B_\Theta=\Bbbk[x,\eta]$, so
\[
\mathfrak{G}_1(B_\Theta)
\;=\;
\frac{\Bbbk[x,\eta]\langle\alpha_1,\beta_1\rangle}
{\langle\alpha_1^2,\;\beta_1^2,\;
\alpha_1\beta_1+\beta_1\alpha_1-\eta^2\rangle}.
\]
As a left $B_\Theta$-module, $\mathfrak{G}_1(B_\Theta)$ is free of
rank~$4$ with basis $\{1,\alpha_1,\beta_1,\alpha_1\beta_1\}$.
The Clifford relation is
\[
\alpha_1\beta_1+\beta_1\alpha_1 \;=\; \eta^2.
\]

\medskip
\noindent\textbf{Step 7: Morita triviality after inverting $u$.}
Set $R:=B_\Theta[u^{-1}]=\Bbbk[x,\eta,\eta^{-2}]$.
By Theorem~\ref{thm:tholog-one-handle-matrix},
\[
\mathfrak{G}_1(R)
\;\cong\;
\underline{\operatorname{End}}_R(R\oplus R\varepsilon),
\qquad |\varepsilon|=-1,
\]
via the explicit isomorphism
\[
\alpha_1\;\longmapsto\;u\cdot(\varepsilon\wedge -),
\qquad
\beta_1\;\longmapsto\;\iota_\varepsilon.
\]
Check the Clifford relation:
\[
u(\varepsilon\wedge -)\,\iota_\varepsilon
+\iota_\varepsilon\,u(\varepsilon\wedge -)
\;=\;
u\,\bigl((\varepsilon\wedge-)\iota_\varepsilon
+\iota_\varepsilon(\varepsilon\wedge-)\bigr)
\;=\;u\cdot\operatorname{id}
\;=\;\eta^2\cdot\operatorname{id}.
\]
This confirms $\mathfrak{G}_1(R)\cong\operatorname{Mat}_{2\times 2}(R)$,
so
\[
D\bigl(\mathfrak{G}_1(R)\text{-}\mathbf{mod}\bigr)
\;\simeq\;
D\bigl(R\text{-}\mathbf{mod}\bigr).
\]
The genus-$1$ completion is derived Morita-trivial after inverting
$u=\eta^2$, confirming
Theorem~\ref{thm:tholog-morita-triviality} in this case.

For general genus~$g$, by
Theorem~\ref{thm:tholog-morita-triviality},
\[
\mathfrak{G}_g(R)
\;\cong\;
\operatorname{Mat}_{2^g\times 2^g}(R),
\]
so all higher-genus data is Morita-trivial away from $u=0$.

\medskip
\noindent\textbf{Step 8: the strict locus $u=0$.}
Modulo $u=\eta^2$, by Theorem~\ref{thm:tholog-exterior-degeneration},
\[
\mathfrak{G}_g(B_\Theta)/(\eta^2)
\;\cong\;
\bigl(\Bbbk[x,\eta]/(\eta^2)\bigr)
\otimes
\Lambda(\alpha_1,\beta_1,\dots,\alpha_g,\beta_g).
\]
Since $\Bbbk[x,\eta]/(\eta^2)=\Bbbk[x]\otimes\Lambda(\eta)$
is an exterior-algebra extension of $\Bbbk[x]$, the strict-locus
genus completion is
\[
\Bbbk[x]\otimes\Lambda(\eta,\alpha_1,\beta_1,\dots,\alpha_g,\beta_g),
\]
an exterior algebra on $2g+1$ odd generators over $\Bbbk[x]$.
This recovers the expected exterior handle algebra for the abelian
theory on the secondary-anomaly cone.

\medskip
\noindent\textbf{Summary.}
For $U(1)$ Chern--Simons at level $k\neq 0$:
\begin{center}
\renewcommand{\arraystretch}{1.2}
\begin{tabular}{ll}
Input algebra & $B=\Bbbk[x]$, $|x|=2$, $d=0$ \\
Anomaly class & $\Theta=kx\in B^2$ (central, closed) \\
Transgression algebra &
 $B_\Theta=\Bbbk[x,\eta]$, $d\eta=kx$ \\
Cohomology &
 $H^\bullet(B_\Theta)\cong\Bbbk[u]$ \\
Secondary anomaly & $u=\eta^2$, closed and central \\
Genus-$g$ Clifford completion &
 $\mathfrak{G}_g(B_\Theta)$: Clifford on $2g$ generators
 with $\{\alpha_i,\beta_j\}=\delta_{ij}\eta^2$ \\
Morita triviality &
 $\mathfrak{G}_g(B_\Theta)[\eta^{-2}]
 \cong\operatorname{Mat}_{2^g}(B_\Theta[\eta^{-2}])$ \\
Strict locus &
 $\mathfrak{G}_g(B_\Theta)/(\eta^2)
 \cong\Bbbk[x]\otimes\Lambda^{2g+1}$
\end{tabular}
\end{center}
Every abstract construction of
Sections~\ref{subsec:tholog-transgression-algebra}--\ref{subsec:tholog-morita-triviality}
is visible and computable by hand.
\end{example}

\subsection{The Virasoro algebra: anomaly completion
at the shadow obstruction tower}
% label removed: subsec:tholog-virasoro

The Virasoro algebra is the simplest $W$-algebra whose
anomaly-completed datum is nontrivial. Unlike the abelian
example, the shadow obstruction tower is \emph{infinite} (type~$M$), the
curvature is a genuine function of the central charge, and
the transgression algebra connects two Virasoro algebras at
complementary central charges.

\begin{computation}[Virasoro anomaly-completed datum;
\ClaimStatusProvedHere]
% label removed: comp:virasoro-anomaly-completion
\index{Virasoro algebra!anomaly completion|textbf}
\index{anomaly completion!Virasoro}
Let $\mathrm{Vir}_c$ denote the Virasoro vertex algebra at
central charge~$c$.

\emph{Koszul duality.}
The chiral Koszul dual is
(Vol~I, Theorem~A):
\begin{equation}% label removed: eq:vir-koszul-dual
\mathrm{Vir}_c^!
\;\simeq\;
\mathrm{Vir}_{26-c}.
\end{equation}
The duality is \emph{same-family}: both sides are Virasoro
algebras, at complementary central charges summing to~$26$.
The self-dual point is $c=13$, where
$\mathrm{Vir}_{13}^!\simeq\mathrm{Vir}_{13}$.

\emph{The curvature as a function of~$c$.}
The modular Koszul curvature $\kappa(\mathrm{Vir}_c)\in\mathbb Q(c)$
is computed from the cyclic pairing on the bar complex of
$\mathrm{Vir}_c$ (Vol~I, Theorem~D). Its key properties:
\begin{enumerate}[label=\textup{(\roman*)}]
\item $\kappa(\mathrm{Vir}_{13})=\tfrac{13}{2}$ (self-duality
 at $c=13$ fixes the curvature to the half-complementarity value);
\item $\kappa(\mathrm{Vir}_c)\neq 0$ for $c\neq 13$
 generically;
\item $\kappa$ is duality-constrained:
 $\kappa(\mathrm{Vir}_c)+\kappa(\mathrm{Vir}_{26-c})=13$
 (Theorem~D; for Virasoro, the complementarity sum is a nonzero constant).
\end{enumerate}
The curvature measures the \emph{distance from self-duality}:
it is the scalar obstruction to extending the genus-$0$
Koszul duality to genus~$1$.

\emph{The transgression algebra as bridge.}
Set $B=\bar B^{\mathrm{ch}}(\mathrm{Vir}_c)$
(the bar coalgebra as dg~algebra). The anomaly class is
$\Theta=\kappa(\mathrm{Vir}_c)\cdot\omega_1$, and the
transgression algebra
\[
B_\Theta
\;=\;
\frac{B*\Bbbk\langle\eta\rangle}
{\langle\eta b-(-1)^{|b|}b\eta\rangle},
\qquad
d\eta=\kappa\cdot\omega_1,
\]
is the \emph{universal bridge} between $\mathrm{Vir}_c$
and $\mathrm{Vir}_{26-c}$. A $B_\Theta$-module $M$
carries simultaneously:
\begin{itemize}
\item a $\mathrm{Vir}_c$-module structure (via the
 $B$-action);
\item a chain-level homotopy $h_M=\eta_M$ relating the
 $\mathrm{Vir}_c$-action to the $\mathrm{Vir}_{26-c}$-action
 (via the bar-cohomology/Koszul structure).
\end{itemize}
The neutralisation condition $dh_M=\Theta_M$ is the statement
that the genus-$1$ curvature obstruction for $M$ can be
absorbed into a homotopy connecting the two central charges.

\emph{The secondary anomaly and the shadow obstruction tower.}
The secondary anomaly $u=\eta^2$ is closed and central in
$B_\Theta$. The Virasoro shadow obstruction tower from Vol~I
(type~$M$, infinite depth) produces nonvanishing shadow
projections $\mathrm{Sh}_r(\Theta_{\mathrm{Vir}})$ at every
degree $r\ge 2$: the curvature~$\kappa$ at degree~$2$,
the cubic shadow~$C$ at degree~$3$, the quartic resonance
class $Q^{\mathrm{contact}}_{\mathrm{Vir}}
=10/[c(5c{+}22)]$ at degree~$4$, and the quintic obstruction
(forced nonvanishing, degree~$5$).
The transgression algebra $B_\Theta$ encodes
\emph{all} of these shadows simultaneously: the secondary
anomaly $u=\eta^2$ is the genus-raising operation, and the
shadow projections are the $u$-adic Taylor coefficients of
the full MC element $\Theta_{\mathrm{Vir}}$.

\emph{Genus structure.}
The genus-$g$ Clifford completion
$\mathfrak G_g(B_\Theta)$ has the standard form:
\begin{itemize}
\item Away from $u$-locus ($\eta^2$ invertible):
 Morita-trivial, $\mathfrak G_g[\eta^{-2}]\cong
 \operatorname{Mat}_{2^g}$. The $\mathrm{Vir}_c$- and
 $\mathrm{Vir}_{26-c}$-module categories are
 interchangeable at every genus.
\item On the strict locus ($\eta^2=0$):
 the curvature vanishes on the strict quadratic locus
 ($c=0$ for $\mathrm{Vir}_c$, or $c=26$ for
 $\mathrm{Vir}_{26-c}$).
 The genus-$g$ extension degenerates to an exterior algebra
 $\Lambda^{2g+1}\otimes B/(\Theta)$.
 The Koszul self-dual point $c=13$ is distinct:
 $\kappa = 13/2 \neq 0$ and $\eta^2 \neq 0$.
\end{itemize}

\emph{Why the Virasoro is the canonical example.}
The Virasoro is the unique vertex algebra where:
\begin{enumerate}[label=\textup{(\arabic*)}]
\item the Koszul dual is the \emph{same type} of algebra
 (same-family duality);
\item the curvature is a \emph{single rational function}
 of one parameter ($c$);
\item the shadow obstruction tower is \emph{infinite} (type~$M$),
 showing that the transgression algebra carries genuinely
 infinite-depth data;
\item the self-dual point ($c=13$) is an isolated point in
 moduli, not a special sublocus.
\end{enumerate}
Every structural feature of the anomaly-completed machine (the
transgression bridge, the secondary anomaly, the Morita
triviality, the strict-locus degeneration) is visible and
controlled by the single parameter~$c$.
\end{computation}

\subsection{The $\mathfrak{sl}_3$ DS chain:
anomaly-completed holography from the nilpotent poset}
% label removed: subsec:tholog-ds-chain

The Virasoro example is the principal reduction of
$\widehat{\mathfrak{sl}}_2$. The first example where the
DS hierarchy and the anomaly-completed machine interact
nontrivially is the $\mathfrak{sl}_3$ chain:
\[
\widehat{\mathfrak{sl}}_{3,k}
\;\xrightarrow{\;\mathrm{DS}_{\min}\;}
\mathcal B^k
\;\xrightarrow{\;\mathrm{DS}_{\mathrm{sub}\to\mathrm{prin}}\;}
\mathcal W_3.
\]
Each arrow is a quantum Drinfeld--Sokolov reduction, and each
vertex carries its own anomaly-completed datum. The key
structural insight is that DS reduction is a \emph{functor on
anomaly-completed data}: it transforms the transgression
algebra, the curvature, and the genus-Clifford completion
coherently.

\begin{construction}[The $\mathfrak{sl}_3$ chain as
anomaly-completed diagram; \ClaimStatusProvedHere]
% label removed: constr:sl3-anomaly-chain
\index{Drinfeld--Sokolov reduction!anomaly-completed functor}
\index{sl3 DS chain@$\mathfrak{sl}_3$ DS chain|textbf}
At each vertex of the nilpotent poset of $\mathfrak{sl}_3$
(Vol~I, Computation~\ref*{V1-comp:sl3-ds-hierarchy}), the
anomaly-completed datum
$(\mathcal A,\;\mathcal A^!,\;\kappa,\;\Theta,\;B_\Theta)$
is:
\begin{center}
\renewcommand{\arraystretch}{1.3}
\begin{tabular}{lccc}
& \textbf{Level~$0$} & \textbf{Level~$1$} & \textbf{Level~$2$} \\
\hline
Orbit & $\{0\}$ & $(2,1)$ & $(3)$ \\[2pt]
$\mathcal A$ &
 $\widehat{\mathfrak{sl}}_{3,k}$ &
 $\mathcal B^k$ &
 $\mathcal W_3^k$ \\[2pt]
$\mathcal A^!$ &
 $\widehat{\mathfrak{sl}}_{3,-k-6}$ &
 $\mathcal B^{-k-6}$ &
 $\mathcal W_3^{-k-6}$ \\[2pt]
$c+c'$ & $16$ & $196$ & $100$ \\[2pt]
Ghost dim & $0$ & $2$ & $6$ \\[2pt]
$c_{\mathrm{ghost}}$ &
 $0$ &
 $\frac{2k-3}{k+3}$ &
 $6$ \\[2pt]
$\kappa=0$ at & $k=-3$ (crit.) &
 $c=98$ & $c=50$ \\[2pt]
Dual orbit & $\{0\}$ & $(2,1)$ & $(3)$ \\[2pt]
Cross-orbit? & no & no (self-$t$) & no
\end{tabular}
\end{center}

\emph{The vertical structure.}
At each level, the anomaly class $\Theta_i=\kappa_i\cdot\omega_1$
is closed and central in the bar dg~algebra $B_i$, and the
transgression algebra $(B_i)_{\Theta_i}$ bridges the
$W$-algebra $\mathcal A_i$ with its Koszul dual
$\mathcal A_i^!$.

\emph{The horizontal structure: DS as functor.}
The DS reduction $\mathrm{DS}_f\colon\mathcal A\to\mathcal W$
extends to the anomaly-completed level. The bar complex
transforms as
$B_{\mathcal W}\simeq H^0_{\mathrm{BRST}}(B_{\mathcal A}
\otimes B_{\mathrm{ghost}})$,
and the curvature transforms as
\begin{equation}% label removed: eq:curvature-ds-transform
\kappa(\mathcal W)
\;=\;
\kappa(\mathcal A)
\;-\;
\kappa_{\mathrm{ghost}}(f,k),
\end{equation}
where $\kappa_{\mathrm{ghost}}(f,k)$ is the ghost system's
contribution to the modular characteristic. For the principal
reduction of $\mathfrak{sl}_N$, $\kappa_{\mathrm{ghost}}$ is
level-independent; for non-principal reductions,
$\kappa_{\mathrm{ghost}}$ depends on~$k$, which is the source
of the level-dependent ghost central charge
$c_{\mathrm{ghost}}=(2k{-}3)/(k{+}3)$ at level~$1$.

\emph{Cross-orbit duality from the chain.}
At level~$1$, the partition $(2,1)$ is self-transpose, so
the Koszul duality is same-orbit:
$(\mathcal B^k)^!\simeq\mathcal B^{-k-6}$. The
anomaly-completed datum at level~$1$ is obtained from the
datum at level~$0$ by applying
$\mathrm{DS}_{\min}$: the curvature shifts by
$\kappa_{\mathrm{ghost}}$, the transgression algebra
acquires the ghost sector, and the genus-Clifford completion
inherits a ghost-extended handle algebra. The
\emph{cross-orbit phenomenon} for other nilpotent types arises
when the DS reduction changes the orbit type of the Koszul
dual: the functoriality of DS on anomaly-completed data
then forces the transgression algebra to bridge two
\emph{different} $W$-algebra families.
\end{construction}

\begin{theorem}[DS functoriality on anomaly-completed data;
\ClaimStatusProvedHere]
% label removed: thm:ds-anomaly-functoriality
\index{Drinfeld--Sokolov reduction!functoriality on anomaly completion}
Let $\mathcal A$ be a vertex algebra with anomaly-completed
datum $(B,\Theta,B_\Theta)$ and let $f\in\mathfrak g$ be a
nilpotent element with DS reduction
$\mathcal W=\mathrm{DS}_f(\mathcal A)$. Then:
\begin{enumerate}[label=\textup{(\roman*)}]
\item The bar complex transforms:
 $B_{\mathcal W}\simeq H^0_{\mathrm{BRST}}
 (B_{\mathcal A}\otimes\mathcal F_{\mathrm{gh}})$.
\item The curvature transforms by~\eqref{eq:curvature-ds-transform}.
\item The transgression algebra transforms:
 $(B_{\mathcal W})_{\Theta_{\mathcal W}}
 \simeq
 H^0_{\mathrm{BRST}}\bigl(
 (B_{\mathcal A})_{\Theta_{\mathcal A}}
 \otimes\mathcal F_{\mathrm{gh}}
 \bigr)$.
\item The genus-$g$ Clifford completion transforms:
 $\mathfrak G_g(B_{\mathcal W})
 \simeq
 H^0_{\mathrm{BRST}}\bigl(
 \mathfrak G_g(B_{\mathcal A})
 \otimes\mathcal F_{\mathrm{gh}}
 \bigr)$.
\end{enumerate}
In particular, if the anomaly-completed datum at
$\mathcal A$ is Morita-trivial at genus~$g$ (away from the
$u$-locus), then so is the datum at $\mathcal W$.
\end{theorem}

\begin{proof}
The DS functor $\mathrm{DS}_f=H^0((-)\otimes
\mathcal F_{\mathrm{gh}},Q_{\mathrm{DS}})$ is exact on the
category of vertex algebras with bounded-below weight grading,
and commutes with the bar construction by the BRST--bar
intertwining of Vol~I (Theorem~A applied to the BRST
extension). The transgression algebra is determined by
$(B,\Theta)$, and both transform functorially under BRST
cohomology: the generator $\eta$ survives BRST reduction
(it commutes with $Q_{\mathrm{DS}}$ because $\Theta$ is
central), and $d\eta=\Theta$ transforms to
$d\eta=\Theta_{\mathcal W}$ by~(ii). The Clifford generators
$\alpha_i,\beta_j$ likewise commute with $Q_{\mathrm{DS}}$
(they are purely topological, living in the genus factor),
so~(iv) follows from~(iii).
\end{proof}

\begin{remark}[The $\mathfrak{sl}_3$ chain as the
platonic model]
% label removed: rem:sl3-platonic-model
The $\mathfrak{sl}_3$ DS chain is the simplest instance of
the general principle: the nilpotent poset of~$\mathfrak g$
indexes a \emph{diagram} of anomaly-completed data, with the
DS functors as the morphisms. The anomaly class $\Theta_i$ at
each vertex is \emph{not} independent: it is determined by the
curvature transform~\eqref{eq:curvature-ds-transform} from
the affine-level datum. The self-dual locus ($\kappa=0$)
varies across the diagram: at level~$0$ it is $k=-3$
(critical), at level~$1$ it is $c=98$, and at level~$2$ it
is $c=50$. The anomaly-completed machine at each vertex is
the \emph{fibre} of this diagram over the specific nilpotent
orbit; the DS functors are the \emph{transition maps} between
fibres.

For $\mathfrak{sl}_N$ with hook partitions, the same
structure governs the entire reduction network $\Gamma_N$
(Vol~I, Proposition~\ref*{V1-prop:transport-propagation-frontier}):
the transport-propagation lemma is the statement that the
anomaly-completed diagram is \emph{connected}, and the
transport-to-transpose conjecture
(Conjecture~\ref*{V1-conj:type-a-transport-to-transpose-frontier})
is the statement that it is \emph{complete}.
\end{remark}

\subsection{Hook-type W-algebras: cross-orbit anomaly completion}
% label removed: subsec:tholog-hook-type

The abelian Chern--Simons example above is self-dual: the boundary
algebra and its Koszul dual are of the same type. For non-principal
$W$-algebras, the anomaly-completed machine acquires a fundamentally
new feature: the \emph{Koszul dual is a different $W$-algebra},
associated to the transpose partition. The transgression algebra
$B_\Theta$ becomes the bridge between two distinct $W$-algebra families,
and the genus-Clifford completion carries both algebras simultaneously.

\begin{construction}[Hook-type anomaly-completed datum;
\ClaimStatusProvedHere]
% label removed: constr:hook-type-anomaly-datum
\index{W-algebra@$\mathcal W$-algebra!hook-type!anomaly completion|textbf}
\index{anomaly completion!non-principal W-algebra}
Let $\lambda=(N{-}r,1^r)\vdash N$ be a hook partition of~$N$
with $1\le r\le N{-}2$, and let $\mathcal W_\lambda^k
:= \mathcal W^k(\mathfrak{sl}_N,f_\lambda)$ be the corresponding
$W$-algebra at level~$k\neq -N$. The anomaly-completed
holographic datum is assembled in five steps.

\emph{Step~1: The BRST input.}
The quantum Drinfeld--Sokolov reduction produces
$\mathcal W_\lambda^k$ as the BRST cohomology
\[
\mathcal W_\lambda^k
\;=\;
H^0\!\bigl(
\widehat{\mathfrak{sl}}_{N,k}
\otimes
\mathcal F_{\mathrm{gh}}(\mathfrak m^+_\lambda),\;
Q_{\mathrm{DS}}^{(\chi_\lambda)}
\bigr),
\]
where $\mathfrak m^+_\lambda$ is the nilradical of the
Jacobson--Morozov parabolic for $f_\lambda$,
$\dim(\mathfrak m^+_\lambda)=r(N{-}r)$, and $\chi_\lambda$
is the Slodowy character. The $W$-algebra has strong generators
of conformal weights $1,2,\dots,r{+}1$ (from the
$\mathfrak{sl}_2$-grading of
$\mathfrak{sl}_N^{f_\lambda}$).

\emph{Step~2: Cross-orbit Koszul duality.}
By the hook-type transport theorem
(Vol~I, Theorem~\ref*{V1-thm:hook-transport-corridor}),
at generic level:
\begin{equation}% label removed: eq:hook-koszul-duality
(\mathcal W_\lambda^k)^!
\;\simeq\;
\mathcal W_{\lambda^t}^{k^\vee},
\qquad
k^\vee=-k-2N,
\qquad
\lambda^t=(r{+}1,1^{N{-}r{-}1}).
\end{equation}
The partition transpose $\lambda\mapsto\lambda^t$ is the type-A
specialisation of Barbasch--Vogan duality. For $r\neq(N{-}1)/2$,
the partitions $\lambda\neq\lambda^t$ and the Koszul dual is a
\emph{genuinely different} $W$-algebra.

The central-charge duality identity:
\begin{equation}% label removed: eq:hook-cc-sum
c(\mathcal W_\lambda^k)
+c(\mathcal W_{\lambda^t}^{k^\vee})
\;=\;
c_{\mathrm{sum}}(N,r),
\end{equation}
where $c_{\mathrm{sum}}(N,r)$ is the Freudenthal--de~Vries
constant for the pair $(\lambda,\lambda^t)$. For the
Bershadsky--Polyakov algebra ($N=3$, $r=1$):
$c_{\mathrm{sum}}(3,1)=196$
(Vol~I, Conjecture~\ref*{V1-conj:bp-duality}).

\emph{Step~3: The curvature and anomaly class.}
The chiral Koszul curvature $\kappa(\mathcal W_\lambda^k)$ is
the modular characteristic of $\mathcal W_\lambda^k$,
computed from the bar complex via the cyclic pairing
(Vol~I, Theorem~D). It is a rational function of the
level~$k$ that vanishes at the self-dual point
$c(\mathcal W_\lambda^k)=\tfrac{1}{2}\,c_{\mathrm{sum}}$.
For generic $k$,
$\kappa\neq 0$ and the genus-$1$ curvature is
$d^2=\kappa\cdot\omega_1$. The anomaly class is
$\Theta=\kappa\cdot\omega_1\in Z^2(B)\cap Z(B)$, where $B$
is the bar dg~algebra of the boundary system.

\emph{Step~4: The transgression algebra.}
Set $B=\bar B^{\mathrm{ch}}(\mathcal W_\lambda^k)$
(the bar coalgebra viewed as dg~algebra via the convolution
product). One expects the transgression algebra to be
\[
B_\Theta
\;=\;
\frac{B*\Bbbk\langle\eta\rangle}
{\langle\eta b-(-1)^{|b|}b\eta\rangle},
\qquad
d\eta=\Theta=\kappa\cdot\omega_1,
\quad|\eta|=1.
\]
The secondary anomaly is then expected to be
$u=\eta^2\in Z^2(B_\Theta)\cap Z(B_\Theta)$.

The new feature expected for non-principal $W$-algebras is that the
$B_\Theta$-modules carry data from \emph{both} $\mathcal W_\lambda^k$
and its Koszul dual $\mathcal W_{\lambda^t}^{k^\vee}$. A
$\Theta$-neutralisation
$h_M$ with $dh_M=\Theta_M$ on a $B$-module $M$ is an intertwining
datum between the $\mathcal W_\lambda^k$-action (via $B$) and
the $\mathcal W_{\lambda^t}^{k^\vee}$-action (via the Koszul
dual structure on bar cohomology).

\emph{Step~5: Genus-Clifford completion.}
One then expects the genus-$g$ Clifford completion
$\mathfrak G_g(B_\Theta)$ to have $2g$ generators
$\alpha_1,\beta_1,\dots,\alpha_g,\beta_g$ with
$\alpha_i\beta_j+\beta_j\alpha_i=\delta_{ij}\,\eta^2$ and
$\alpha_i^2=\beta_j^2=0$.

\begin{enumerate}[label=\textup{(\roman*)}]
\item \emph{Away from the $u$-locus} ($\eta^2$ invertible):
 $\mathfrak G_g(B_\Theta)[\eta^{-2}]
 \cong\operatorname{Mat}_{2^g}(B_\Theta[\eta^{-2}])$
 (Morita triviality).
 Higher genus is then expected to be invisible: the $\mathcal W_\lambda^k$-module
 structure at genus~$g$ is expected to be Morita-equivalent to genus~$0$.
 This is the regime where the cross-orbit duality is expected to be
 ``unobstructed'': both $\mathcal W_\lambda^k$ and
 $\mathcal W_{\lambda^t}^{k^\vee}$ are expected to contribute to the
 full module category without interference.

\item \emph{On the strict locus} ($u=\eta^2=0$):
 $\mathfrak G_g(B_\Theta)/(\eta^2)
 \cong B/(\Theta)\otimes\Lambda^{2g+1}$.
 The higher-genus structure is expected to degenerate to an exterior-algebra
 extension. On this locus, the curvature is expected to vanish
 ($\kappa=0$) and the cross-orbit duality is expected to collapse: the
 algebra and its dual are expected to become indistinguishable at the level
 of the anomaly-completed module category.
\end{enumerate}
\end{construction}

\begin{computation}[Bershadsky--Polyakov algebra:
$\mathfrak{sl}_3$, hook $(2,1)$; \ClaimStatusConditional]
% label removed: comp:bp-anomaly-completion
\index{Bershadsky--Polyakov algebra!anomaly completion|textbf}
\index{W-algebra@$\mathcal W$-algebra!subregular!anomaly completion}
The Bershadsky--Polyakov algebra
$\mathcal B^k=\mathcal W^k(\mathfrak{sl}_3,f_{(2,1)})$
is the simplest non-principal case, with $N=3$, $r=1$.

\emph{Generators.}
Five strong generators: the $\widehat{\mathfrak{sl}}_2$
subalgebra $\{e,f,h\}$ (conformal weight~$1$) and two
fermionic generators $G^\pm$ (conformal weight~$3/2$).
Equivalently, this is the $\mathcal N=2$ superconformal
algebra at $c=3(2k{+}1)/(k{+}3)$.

\emph{Koszul dual.}
The partition $(2,1)$ is self-transpose:
$(2,1)^t=(2,1)$. Therefore:
\[
(\mathcal B^k)^!
\;\simeq\;
\mathcal B^{-k-6},
\]
and the Koszul dual is the \emph{same type} of $W$-algebra
at the Feigin--Frenkel dual level $k^\vee=-k-6$.
Conditional on the BP duality conjecture of Volume~I,
\[
c(\mathcal B^k)+c(\mathcal B^{-k-6})=196.
\]

\emph{Curvature.}
\[
\kappa(\mathcal B^k)
\;\in\;\mathbb Q(k),
\]
a rational function of~$k$ computed from the cyclic pairing
on the bar complex of $\mathcal B^k$ (Vol~I, Theorem~D).
Since $(\mathcal B^k)^!=\mathcal B^{-k-6}$ and the
conjectural central-charge sum is $c+c'=196$, the predicted
self-dual point is $c=98$ and the curvature is generically
nonzero.
The explicit formula for $\kappa(\mathcal B^k)$ is determined
by the bar complex computation
(Computation~\ref*{V1-comp:sl3-ds-hierarchy} in Vol~I):
it receives contributions from both the
$\widehat{\mathfrak{sl}}_2$ sector (double poles) and the
fermionic $G^\pm$ sector (triple pole in the $G^+G^-$ OPE).

\emph{Anomaly class and transgression.}
Within this conditional package, the anomaly class is
$\Theta=\kappa\cdot\omega_1\in Z^2(B)$,
where $B=\bar B^{\mathrm{ch}}(\mathcal B^k)$. The
transgression algebra $B_\Theta$ carries the combined
structure of $\mathcal B^k$ and $\mathcal B^{-k-6}$:
the neutralisation datum $h_M$ intertwines between the two
levels.

\emph{Ghost system.}
The DS reduction from $\widehat{\mathfrak{sl}}_3$ to
$\mathcal B^k$ removes $r(N{-}r)=1\cdot 2=2$ positive roots,
producing a ghost system with central charge
$c_{\mathrm{ghost}}=(2k{-}3)/(k{+}3)$
(level-dependent, unlike the principal case where
$c_{\mathrm{ghost}}=N(N{-}1)$ is constant). The level
dependence of $c_{\mathrm{ghost}}$ is reflected in the
anomaly class: the ghost contribution to~$\Theta$ shifts
the self-dual level away from $k=-h^\vee/2$.

\emph{Secondary anomaly and genus structure.}
Within the same conditional package,
$u=\eta^2$ is closed and central in $B_\Theta$. The
genus-$1$ Clifford completion
$\mathfrak G_1(B_\Theta)=B_\Theta\langle\alpha_1,\beta_1
\mid\alpha_1\beta_1+\beta_1\alpha_1=\eta^2\rangle$
governs the torus partition function of the anomaly-completed
$\mathcal B^k$ theory. After inverting $\eta^2$:
the $\mathcal B^k$-modules and $\mathcal B^{-k-6}$-modules
become Morita-equivalent (genus invisible). On the strict
locus $\eta^2=0$: the two families decouple and the genus
structure degenerates.
\end{computation}

\begin{remark}[Mukai-doubling Koszul conductor $K^{\kappa_{\mathrm{ch}}}=8$ on the $\mathcal B$-family; \ClaimStatusConjectured]
\label{rem:anomtop-mukai-koszul-conductor-8}
\index{Koszul conductor!Mukai-doubling!B-family}
\index{Bruinier Heegner reciprocity!three faces}
The Bershadsky--Polyakov $\mathcal B$-family above carries a
single numerical invariant that refines the transgression
theorem to a Bruinier-Heegner-reciprocal statement. Let
$\kappa_{\mathrm{ch}}$ denote the curvature class
$\kappa(\mathcal B^{k})$ entering $\Theta$. Then the
\emph{Koszul conductor} of the $\mathcal B$-family under the
Mukai-doubling (of the ambient lattice on which the
Vol~III CY-to-chiral functor $\Phi$ is evaluated) is
\[
K^{\kappa_{\mathrm{ch}}}
\;=\;
8,
\]
the unique value for which the three faces of the
$\mathcal B$-family coincide under Bruinier Heegner
reciprocity:
\begin{enumerate}[label=\textup{(\roman*)}]
\item the \emph{chiral face}: the Koszul-curvature obstruction to
extending $\mathcal B^{k}\leftrightarrow\mathcal B^{-k-6}$ from
genus~$0$ to genus~$1$;
\item the \emph{transgression face}: the square
$u=\eta^{2}$ of the odd generator adjoined in
Definition~\ref{def:tholog-transgression-algebra};
\item the \emph{Heegner face}: the level of the Bruinier
Heegner divisor on $\overline{\mathcal A_{2}}$ at which the
paramodular cocycle of
Remark~\ref{rem:anomtop-delta5-paramodular-output}
acquires its first nontrivial residue.
\end{enumerate}
Under the universal $\Phi$-trace identity of
Remark~\ref{rem:anomtop-phi-trace-k3xt2}, the three faces
are equal because all three read off the same Koszul
conductor $K^{\kappa_{\mathrm{ch}}}=8$ from the
Mukai-doubled lattice. The $\mathcal B$-family is thus the
smallest $W$-algebra family where the transgression theorem,
the chiral Koszul theorem, and the Bruinier Heegner
reciprocity all measure the same integer; its role in the
anomaly-completed holographic programme is to provide a
falsifiable numerical check that the three faces really do
coincide beyond the principal Virasoro case.
\end{remark}

\begin{conjecture}[$\mathfrak{sl}_4$ hook $(2,1,1)$:
cross-orbit anomaly completion;
\ClaimStatusConjectured]
% label removed: comp:sl4-hook-anomaly-completion
\index{W-algebra@$\mathcal W$-algebra!hook-type!sl4@$\mathfrak{sl}_4$}
For $N=4$, $r=2$, the hook partition $\lambda=(2,1,1)$ has
transpose $\lambda^t=(3,1)$. Since
$\lambda\neq\lambda^t$, the Koszul duality is genuinely
cross-orbit:
\[
\bigl(\mathcal W^k(\mathfrak{sl}_4,f_{(2,1,1)})\bigr)^!
\;\simeq\;
\mathcal W^{-k-8}(\mathfrak{sl}_4,f_{(3,1)}).
\]

\emph{The cross-orbit structure.}
The partition $(2,1,1)$ is the minimal nilpotent of
$\mathfrak{sl}_4$; the partition $(3,1)$ is the subregular
nilpotent. The DS reduction for $(2,1,1)$ removes
$r(N{-}r)=2\cdot 2=4$ roots (ghost system of dimension~$4$);
the DS reduction for $(3,1)$ removes
$r'(N{-}r')=1\cdot 3=3$ roots (ghost system of dimension~$3$).
The ghost systems are \emph{different}, reflecting the
geometric asymmetry of the cross-orbit duality.

\emph{The transgression algebra as cross-orbit bridge.}
The transgression algebra $B_\Theta$ is expected to carry a single
differential $d\eta=\Theta=\kappa\cdot\omega_1$ that is expected to
simultaneously encode:
\begin{enumerate}[label=\textup{(\roman*)}]
\item the genus-$1$ obstruction for
 $\mathcal W^k(\mathfrak{sl}_4,f_{(2,1,1)})$-modules
 (via the $B$-action);
\item the genus-$1$ obstruction for
 $\mathcal W^{-k-8}(\mathfrak{sl}_4,f_{(3,1)})$-modules
 (via the bar-cohomology/Koszul dual structure).
\end{enumerate}
The neutralisation $h_M$ is then expected to be the datum that resolves both
obstructions simultaneously: it would be a chain-level witness
that the curvature $\kappa$ can be absorbed into a
homotopy connecting the minimal and subregular
$W$-algebra actions.

\emph{Genus-Clifford completion.}
One expects $\mathfrak G_g(B_\Theta)$ at genus~$g$ to carry handle
operators $\alpha_i,\beta_i$ with
$\{\alpha_i,\beta_j\}=\delta_{ij}\eta^2$. Since
$\lambda\neq\lambda^t$, the strict locus $\eta^2=0$ is expected to be
the locus where the cross-orbit duality degenerates:
on this locus, the minimal and subregular $W$-algebras
are expected to decouple completely. Away from the strict locus, the
Morita triviality
$\mathfrak G_g[\eta^{-2}]\cong\operatorname{Mat}_{2^g}$
is expected to mean that at every genus, the minimal and subregular
$W$-algebra modules are woven into a single indecomposable
package.
\end{conjecture}

\begin{remark}[The general hook-type pattern; \ClaimStatusConjectured]
% label removed: rem:hook-anomaly-general
\index{anomaly completion!hook-type pattern}
For the full hook family
$\mathcal W^k(\mathfrak{sl}_N,f_{(N-r,1^r)})$:
\begin{center}
\renewcommand{\arraystretch}{1.2}
\begin{tabular}{lll}
\textbf{Datum} & \textbf{Formula} & \textbf{Type} \\
\hline
Partition & $\lambda=(N{-}r,1^r)$ & hook \\
Transpose & $\lambda^t=(r{+}1,1^{N{-}r{-}1})$ & hook \\
Self-dual? & iff $N{-}r=r{+}1$, i.e.\ $N=2r{+}1$ (odd) & \\
Dual level & $k^\vee=-k-2N$ & Feigin--Frenkel \\
Ghost dim & $r(N{-}r)$ & nilradical \\
Ghost dim$'$ & $(N{-}r{-}1)(r{+}1)$ & dual nilradical \\
Cross-orbit? & iff $r\neq(N{-}1)/2$ & genuine \\
\end{tabular}
\end{center}
The expected cross-orbit anomaly completion would be the geometrical
expression of the fact that Koszul duality for non-principal
$W$-algebras does not preserve the nilpotent type: the
anomaly $\Theta$ is expected to be the obstruction to extending the
genus-$0$ cross-orbit duality to higher genus, and the
transgression algebra $B_\Theta$ is expected to be the universal resolution
of this obstruction.
\end{remark}

\section{Synthesis}
\label{sec:acth-supplement}

\begin{remark}[Pentagon anomaly at $\hbar^3$ in topological holography]
\label{rem:acth-1}
Anomaly-completed topological holography at the K3 base
point carries the pentagon coboundary at $\hbar^3$:
\[
 \phi^{(3)}
 \;=\;
 \zeta(3)\, c_{\mathrm{symm}}
 \;+\;
 \tfrac{25}{3}\, c_{\mathrm{timelike}}
 \;+\;
 \tfrac{\Phi_{10}^{\mathrm{un}}}{\eta^{24}}\, c_{\Phi_{10}^{\mathrm{un}}},
\]
with three linearly independent third-order closure pieces. The
$(25/3)$-piece is the timelike-anomaly matching $c_{2d} = -214 = -2\cdot 107$,
read as the Chacaltana--Distler first-principles class-$\mathcal S$
value $c_{4d}(A_1, \Sigma_{0,24}) = 107/6$ via the Beem--Rastelli
factor $-12$. The pentagon constant $25/3$ arises from the $25\cdot 2
= 50$ orthogonal decomposition $c_{2d} = -2\cdot 107$ normalised by the
triple-leg antisymmetrisation denominator~$3$, \emph{not} from any
naive divisibility of $-214$ by $13$: the prime $107$ enters by
$c_{4d} = (5\cdot 24 - 13)/6 = 107/6$ (Shapere--Tachikawa 2008 \S 3),
and the factor $25/3 = 25/3$ remains the Fake-Monster Cartan rank
$\mathrm{rk}(II_{25,1})/3$ minus the timelike-direction subtraction,
independent of $c_{2d}$
\cite{Drinfeld1990,Borcherds1992,GritsenkoNikulin1998}. Cross-ref
Vol.~III Chapter~\ref{V3-ch:k3-chiral-bialgebra-platonic} and
Vol.~I Chapter~\ref{V1-sec:thqg-open-closed-realization} line~1890.
\end{remark}

\begin{remark}[$\widetilde M_{24}$ anomaly as order-$6$ Schur cocycle]
\label{rem:acth-2}
The anomaly-completed topological holography carries the $\widetilde M_{24}$
Schur cocycle $\in H^2(M_{24}, U(1)) = \mathbb Z/12$ of order $6$ as its
topological anomaly, exactly matching the Cheng--Duncan--Harvey 2014
umbral shift $m_\sigma$ on the $M_{24}$-twined K3 elliptic genus. The
Mukai-doubling identity $K^{\kappa_{\mathrm{ch}}} = 8 = 2\cdot 4$
forces the order-$6$ cocycle to lift to $\mathbb Z/12$ under the
$\mathrm{Sym}_5\subset M_{24}$-restriction \cite{ChengDuncanHarvey2014,EguchiOoguriTachikawa2010}.
Cross-ref Chapter~\ref{V3-ch:k3-chiral-bialgebra-platonic}.
\end{remark}

\section{Anomaly polynomial on $\mathrm K3\times\Sigma_{0,24}$ at $c_{2d} = -214$}
\label{sec:acth-supplement-ii}

\begin{remark}[Class-$\mathcal S$ central charge at the K3 point]
\label{rem:anomtop-retract-312}
The Shapere--Tachikawa 2008 class-$\mathcal S$ normalisation
$c_{4d} = (2n_v + n_h)/12$ applied to the Chacaltana--Distler
class-$\mathcal S$ count $(n_v, n_h) = (63, 88)$ for
$(A_1, \Sigma_{0,24})$ gives
\[
c_{4d}=214/12=107/6.
\]
Beem--Rastelli's relation $c_{2d}=-12c_{4d}$ gives
$c_{2d}=-214$. The $6d$ anomaly polynomial, the shadow-tower
spectrum, the mock-modular weight, and the BPS counting below are
read at this central charge.
\end{remark}

The K3 calculation derives $c_{2d} = -214$ from the $6d$ $(2,0)$ $A_1$ anomaly
polynomial integrated over the product geometry, and specifies the
unitary spectrum of class-$\mathcal S$ chiral primaries that the
anomaly labels.

\begin{remark}[$6d$ anomaly-polynomial derivation of $c_{2d}=-214$]
\label{rem:anomtop-anomaly-polynomial}
\index{anomaly polynomial!$6d$ $(2,0)$ $A_1$}
\index{central charge!$c_{2d}=-214$ anomaly-inflow derivation}
The $6d$ $(2,0)$ theory of type $A_1$ on a product geometry
$\mathrm K3 \times \Sigma_{0,24}$ is governed by the anomaly
$8$-form (Intriligator 2000; Alday--Benini--Tachikawa 2009
arXiv:0912.4664, eq.~(2.2)):
\[
 I_8^{A_1}
 \;=\;
 \frac{r_{A_1}}{24}\Bigl(p_2(N) - p_2(T) + \tfrac{1}{4}(p_1(N)-p_1(T))^2\Bigr)
 \;+\;
 \frac{d_{A_1} h^\vee_{A_1}}{24}\,p_2(N),
\]
where $r_{A_1}=1$, $d_{A_1}=3$, $h^\vee_{A_1}=2$ are the rank,
dimension, dual Coxeter number of $\mathfrak{su}(2)$;
$p_i(T)$, $p_i(N)$ are tangent and normal Pontryagin classes.
Integrating over $\mathrm K3$ with
$\int_{\mathrm K3} p_1(T_{\mathrm K3}) = -48$ (Hirzebruch signature
theorem: $\sigma(\mathrm K3) = -16$, $\tau(\mathrm K3) = p_1/3$,
so $\int p_1 = 3\sigma = -48$; equivalently
$\int_{\mathrm K3}\mathrm{td}_2 = 2$ and $\chi(\mathrm K3)=24$)
reduces $I_8$ to a $4$-form $I_4$ on $\Sigma_{0,24}$. Its
coefficient structure contains (Alday--Benini--Tachikawa 2009,
eq.~(2.9)):
\[
 c_{2d} \;=\; \mathrm{coeff}\bigl(p_1(T_\Sigma)\bigr)\cdot
 \bigl[\mathrm{tangent\ contribution}\bigr]
 \;=\; -12\, c_{4d},
\]
the well-known gravitational-anomaly ratio: the $2d$ chiral central
charge equals $-12$ times the $4d$ conformal anomaly $c_{4d}$ of
the corresponding class-$\mathcal S$ theory
(Beem--Lemos--Liendo--Peelaers--Rastelli--van~Rees 2013 eq.~(3.14);
Shapere--Tachikawa 2008 for $c_{4d}$ of $A_1$ on $\Sigma_{0,24}$).
With the correct Chacaltana--Distler 2010 class-$\mathcal S$ count
$(n_v, n_h) = (63, 88)$ on $(A_1, \Sigma_{0,24}, \mathrm{all\ max})$
(eq.~(5.14) therein) and the Shapere--Tachikawa normalisation
$c_{4d} = (2n_v + n_h)/12$,
\[
 c_{4d}(A_1, \Sigma_{0,24})
 \;=\; \frac{2\cdot 63 + 88}{12}
 \;=\; \frac{214}{12}
 \;=\; \frac{107}{6},
\]
we obtain
\[
 c_{2d}(A_1, \Sigma_{0,24})
 \;=\; -12 \cdot \frac{107}{6}
 \;=\; -2 \cdot 107
 \;=\; -214.
\]
The integer $107$ is prime, so $-214 = -2\cdot 107$ admits no
non-trivial integer factorisation apart from $-2\cdot 107$; this is
the 2d chiral anomaly $c_L - c_R$ on the holomorphic half of the
class-$\mathcal S$ chiral algebra, forced by the $6d$ anomaly
polynomial and the orientation of $\Sigma_{0,24}$ relative to the
K3 normal bundle.
\end{remark}

\begin{remark}[Shadow-tower spectrum at $c=-214$;
\ClaimStatusProvedHere]
\label{rem:anomtop-312-spectrum}
\index{spectrum!shadow tower at $c=-214$}
\index{Zamolodchikov!at $c=-214$}
On the stress-tensor line of the class-$\mathcal S$ $A_1$ chiral
algebra at $\Sigma_{0,24}$, the class-$\mathrm M$ shadow tower of
Vol~I Proposition~\ref{V1-prop:stqp-312-factor} specialises to
\begin{align*}
 S_2 \;&=\; \tfrac{c_{2d}}{2} \;=\; -107,
 &
 S_3 \;&=\; 2,
 \\
 S_4 \;&=\; \frac{10}{c(5c+22)} \;=\; \frac{10}{224{,}272}
 \;=\; \frac{5}{112{,}136},
 &
 S_5 \;&=\; \frac{-48}{c^2(5c+22)} \;=\; \frac{3}{3{,}000{,}319},
\end{align*}
all nonzero, with Zamolodchikov denominator
$5c + 22 = 5(-214) + 22 = -1048 \ne 0$ (far from any minimal-model
degeneracy $c_{p,q}$). The quasi-primary norm
$\langle\Lambda|\Lambda\rangle = c(5c+22)/10 = 224{,}272/10
= 22{,}427.2 > 0$, so the weight-$4$ quasi-primary
$\Lambda = {:}TT{:} - (3/10)\partial^2 T$ has positive Zamolodchikov
norm at $c=-214$ -- despite the negative central charge, the norm
sign is $\mathrm{sgn}(c)\cdot\mathrm{sgn}(5c+22) = (-)(-) = +$. The
class-$\mathcal S$ chiral primary spectrum therefore contains at
least the universal Virasoro descendants $\{T, \Lambda, \partial T,
\partial^2 T, \ldots\}$ on the stress-tensor line, labelled by the
shadow tower $(S_2, S_3, S_4, S_5)$ and extending to
depth~$\infty$. The full primary count matches the $24$
$M_{24}$-twined Jacobi forms of Eguchi--Ooguri--Tachikawa 2010:
each $\mathfrak{su}(2)$ puncture contributes an
$\widehat{\mathfrak{sl}}_2$ current at Beem--Rastelli level
$k_{2d} = -N = -2$ per puncture (BLLPRvR 2013 Table~1; Vol~I
Chapter~\ref{V1-sec:thqg-open-closed-realization}),
so the total flavour current algebra is
$\widehat{\mathfrak{sl}}_2^{24}$ at level $k_{2d} = -2$ per factor,
with Sugawara contribution
$c_{\mathrm{Sug}}^{\mathrm{flavour}} = 24\cdot 3\cdot(-2)/(-2+2)$
formally singular at $k = -h^\vee = -2$ (critical level), which is
the correct signal that each $\mathfrak{su}(2)$ flavour factor
sits at its own critical level on the chiral side: the Sugawara
stress tensor is replaced by the Segal--Sugawara centre action, and
the residual Virasoro piece
$c_{2d} - c_{\mathrm{Sug}}^{\mathrm{flavour,\ reg}}$ is read off
from the finite-$k$ Kac--Moody-like construction at critical level
using the reduced Sugawara cocycle of Arakawa 2017
(\texttt{arXiv:1701.05288}) rather than the naive Sugawara formula.
The full primary count at spin~$\le 1$ is
$\dim\{\mathrm{stress\ tensor}\}
+ 24\cdot\dim\{\mathfrak{sl}_2\text{-current}\}
= 1 + 72 = 73$, plus $\Lambda$ and the extended $M_{24}$-twined
tower at higher weight.
\end{remark}

\begin{remark}[$T^2$ modular covariance: Siegel weight $5$, not
elliptic $-c/24$; \ClaimStatusConjectured]
\label{rem:anomtop-T2-modular}
\index{modular covariance!$T^2$ partition function!K3 BKM}
\index{mock modular!Siegel weight 5!Gritsenko--Nikulin}
The naive weight reading $-c_{2d}/24 = 214/24 = 107/12$ is
\emph{not} integer, so the chiral partition function on $T^2$
cannot be packaged as a weight-$-c/24$ elliptic modular form on
$\mathrm{SL}_2(\mathbb Z)$. The correct modular-weight assignment
is Siegel rather than elliptic: the
$\mathrm T^2$ chiral partition function is a Siegel modular form
of weight~$5$ on $\mathrm{Sp}_4(\mathbb Z)$ (or its double cover
carrying the half-integral Gritsenko-lift character), identified
with $\Delta_5 = (\Phi_{10}^{\mathrm{un}})^{1/2}$ through the Borcherds denominator
identity
$\prod_{r\cdot\alpha>0}(1-q^{r\alpha})^{c(r\alpha)} = \Phi_{10}^{\mathrm{un}}$
on $\mathrm{Sp}_4(\mathbb Z)\backslash\mathbb H_2$ (Borcherds 1992
Theorem~10.1; Gritsenko--Nikulin 1998 \S 4). The Siegel weight
$w_{\mathrm{Sieg}} = 5$ is determined by the Borcherds denominator
degree $10$ and the square-root refinement
$\Delta_5 = \sqrt{\Phi_{10}^{\mathrm{un}}}$; it is independent of the elliptic
coordinate $c_{2d}$. The bridge from Borcherds--Siegel weight to
elliptic-fibre weight is the Gritsenko--Nikulin restriction to the
Humbert divisor $H_1$, under which $\Delta_5$ descends to a
weight-$5$ Jacobi form; this is distinct from the naive
Cardy-regime expression $-c_{2d}/24$. Mock completion follows
Zwegers 2002, Cheng--Duncan--Harvey 2014, and
Bringmann--Creutzig--Rolen 2014. Full $T^2$-modular closure on
the holomorphic chiral partition function is conjectural beyond
the $M_{24}$-averaged/$\Phi_{10}^{\mathrm{un}}$-denominator locus; the
gravitational-anomaly form obstructs exact closure (mock shadows
supply the compensator), consistent with Harvey--Moore 1996
threshold integrals. Cross-ref Vol~I
Remark~\ref{V1-rem:stqp-mock-weight}. The non-integrality of
$107/12$ is thus a feature of the $-214$ arithmetic, signalling
Siegel-rather-than-elliptic modular covariance at the K3 base
point.
\end{remark}

\begin{remark}[Effective conformal anomaly and BPS counting at
$c=-214$; \ClaimStatusConditional]
\label{rem:anomtop-ceff-bps}
\index{effective central charge!class-$\mathcal S$ $A_1$ $\Sigma_{0,24}$}
\index{BPS counting!K3 BKM}
The naive Cardy density $\rho(E) \sim e^{2\pi\sqrt{c E/6}}$ is
ill-defined at $c = -214 < 0$. The physically correct counting
function uses the effective central charge
$c_{\mathrm{eff}} = c - 24\, h_{\min}$, with $h_{\min}$ the lowest
eigenvalue of $L_0$ on the BPS module. For $\mathbf H_{\Delta_5}$
the polar part of $1/\Phi_{10}^{\mathrm{un}}(Z)$ at the maximal-rank Siegel cusp
is weight $-10$, so $h_{\min} = -2$ (weight $-10$ divided by
the fugacity-doubling at the separating node gives $-2$ after
Gritsenko--Nikulin 1997 normalisation), and
\[
 c_{\mathrm{eff}} \;=\; -214 + 48 \;=\; -166,
\]
still negative. The BPS index then reads off from the polar part
of the Borcherds denominator: BPS counts match the Harvey--Moore
1996 threshold BPS index, and the sign of $c_{\mathrm{eff}}$
corresponds to counting against the anti-holomorphic orientation
of $\Sigma_{0,24}$ in the $6d$ gravitational anomaly. The Cardy
density is recovered on the full modulus-squared
$|Z_{\mathbf H_{\Delta_5}}|^2$ with Dijkgraaf--Verlinde--Verlinde
1997 evaluation. Conditional on full modular closure of
Remark~\ref{rem:anomtop-T2-modular}: the BPS counting
closes. The effective value $c_{\mathrm{eff}} = -166$ factors as
$-2\cdot 83$ with $83$ prime, mirroring the factorisation
structure of $c_{2d} = -214 = -2\cdot 107$: the shift by
$24 h_{\min} = -48$ moves the prime $107$ to the prime $83$
without disturbing the $2$-adic $2$-factor.
\end{remark}
